\sethyphenation{kannada}{
ಅಂಕು-ಡೊಂಕಾದ
ಅಂಕು-ಶಾಸನ
ಅಂಗ-ಗಳನ್ನು
ಅಂಗ-ಮೋಟ-ನಾಸನ
ಅಂಗ-ವಾಗಿ
ಅಂಗ-ವೈಕಲ್ಯ-ವಿಲ್ಲ-ದವ-ರಾಗಿರ-ಬೇಕು
ಅಂಗುಲ
ಅಂಗುಲ-ಗಳ-ವರೆಗೆ
ಅಂಗುಲ-ಗಳಷ್ಟು
ಅಂಗೈ
ಅಂಗೈ-ಗಳ
ಅಂಗೈ-ಗಳನ್ನು
ಅಂಗೈ-ಗಳ-ವರೆಗೆ
ಅಂಗೈ-ಗಳಿಂದ
ಅಂಗೈ-ಯನ್ನು
ಅಂಟಿ-ಕೊಂಡ-ವರು
ಅಂಟಿ-ನಿಂದ
ಅಂಟು
ಅಂಡ-ಕೋಶ-ಗಳನ್ನು
ಅಂಡ-ಕೋಶದ
ಅಂಡ-ಕೋಶ-ದಿಂದ
ಅಂಡ-ಕೋಶ-ದೊಳಗೆ
ಅಂತರ-ದ-ವರೆಗೆ
ಅಂತಹ
ಅಂದರೆ
ಅಂದಾಜು
ಅಕ್ಕ-ಪಕ್ಕಕ್ಕೆ
ಅಕ್ಕಿ
ಅಕ್ಕಿ-ಯನ್ನು
ಅಕ್ಷರ-ಗಳ
ಅಗಲ
ಅಗಲ-ವಿರ-ಬೇಕು
ಅಗಲ-ವಿರುವ
ಅಗ್ನಿ
ಅಗ್ನಿ-ಹೋತ್ರವು
ಅಚಲ
ಅಚಲ-ವಾಗಿಟ್ಟು-ಕೂಳ್ಳು-ವುದೇ
ಅಚಲ-ವಾಗಿ-ಸು-ವುದೋ
ಅಜಾಸನ
ಅಡಿ-ಯನ್ನು
ಅಡೆ-ತಡೆ-ಗಳನ್ನು
ಅಡೆ-ತಡೆ-ಗಳು
ಅಡ್ಡ-ಲಾಗಿ
ಅಡ್ಡ-ಲಾಗಿ-ರುವ
ಅಡ್ಡ-ಲಾದ
ಅಡ್ಡ-ವಾಗಿ-ರುವ
ಅತಿ
ಅತಿ-ಯಾಗಿ
ಅತಿ-ಯಾದ
ಅತೀ
ಅತ್ಯಂತ
ಅತ್ಯಗತ್ಯ
ಅಥವಾ
ಅದಕ್ಕಿಂತ
ಅದಕ್ಕೆ
ಅದನ್ನು
ಅದರ
ಅದರತ್ತಲೇ
ಅದರಲ್ಲಿಯೂ
ಅದರ-ಲ್ಲಿರುವ
ಅದು
ಅದೂ
ಅದೇ
ಅಧರ್ಮ
ಅಧ್ಯಾಯ
ಅಧ್ಯಾಯವು
ಅನಂತಾಸನ
ಅನಾಹತ
ಅನಿರೀಕ್ಷಿ-ತವಾಗಿ
ಅನುಗ್ರಹಕ್ಕೆ
ಅನುಭವ
ಅನುಭವ-ವಾಗಿ
ಅನುಭವಿ-ಸುತ್ತಾನೋ
ಅನುಮತಿ
ಅನುರಾ-ಗದ
ಅನುಸರಿ-ಸಬೇಕು
ಅನುಸಾರ
ಅನೇಕ
ಅಪಾನ
ಅಪಾರ
ಅಪ್ಪಿ
ಅಪ್ಪಿ-ಕೊಂಡು
ಅಪ್ಪಿ-ಕೊಳ್ಳ-ಬೇಕು
ಅಪ್ಪಿ-ಕೊಳ್ಳುವ
ಅಪ್ರದಕ್ಷಿಣಾ-ಕಾರ-ವಾಗಿ
ಅಭ್ಯಾಸ
ಅಭ್ಯಾಸ-ಕ್ಕಾಗಿ
ಅಭ್ಯಾಸಕ್ಕೆ
ಅಭ್ಯಾಸ-ಗಳಿಂದ
ಅಭ್ಯಾಸ-ಗಳು
ಅಭ್ಯಾಸದ
ಅಭ್ಯಾಸ-ದಲ್ಲಿ
ಅಭ್ಯಾಸ-ದಿಂದ
ಅಭ್ಯಾಸ-ವನ್ನು
ಅಭ್ಯಾಸ-ವಾ-ಗಿದೆ
ಅಭ್ಯಾಸವು
ಅಮರಿ
ಅಮರೋಲಿ
ಅಮರೋಲಿ-ಯನ್ನು
ಅಮೃತ
ಅಮೃತಾ
ಅರ-ಗಿನ
ಅರಗಿ-ನೊಂದಿಗೆ
ಅರಣ್ಯ-ಚಟ-ಕಾಸನ
ಅರಳಿ
ಅರಿಯಲು
ಅರಿ-ವಿರುವ-ವರ
ಅರಿಶಿನ
ಅರೆದು
ಅರ್ಘ್ಯಾಸನ
ಅರ್ಘ್ಯಾಸನ-ದಲ್ಲಿದ್ದಾಗ
ಅರ್ಧ
ಅರ್ಧ-ಚಂದ್ರಾಸನ
ಅರ್ಧ-ದಷ್ಟು
ಅರ್ಧ-ಪಶ್ಚಿಮತಾ-ನಾಸನ
ಅರ್ಹ-ರಲ್ಲ
ಅರ್ಹ-ರಾದ
ಅಲು-ಗಾಡಿಸ-ದೆಯೇ
ಅಲು-ಗಾಡಿಸ-ಬೇಕು
ಅಲ್ಲದೆ
ಅಲ್ಲಿ
ಉಸಿ-ರನ್ನು
ಅಲ್ಲಿಯೇ
ಅಳಲೆ-ಕಾಯಿ
ಅಳಲೆ-ಕಾಯಿ-ಯಿಂದ
ಅಳ-ವಡಿಸಿ
ಅವನ
ಅವ-ನಲ್ಲಿ
ಅವನು
ಅವರ
ಅವ-ರನ್ನು
ಅವರು
ಅವಳ
ಅವಳನ್ನು
ಅವಳಲ್ಲಿ
ಅವಳೊಂದಿಗೆ
ಅವಶ್ಯಕ
ಅವಶ್ಯಕತೆ
ಅವು-ಗಳ
ಅವು-ಗಳನ್ನು
ಅವು-ಗಳು
ಅಶ್ವ-ಸಾಧನಾ
ಅಷ್ಟ-ಕುಂಭಕ
ಅಷ್ಟರ-ಮಟ್ಟಿಗೆ
ಅಷ್ಟಾಂಗ-ಗಳಿಂದ
ಅಷ್ಟೇ
ಅಸಂಖ್ಯಾತ
ಅಸಂಚಯ
ಅಸಮ-ತೋಲನ-ವನ್ನು
ಅಸಲಿ
ಅಸಹನೀಯ-ವಾದ
ಅಹಂಕಾರ
ಅಹಿಂಸೆ
ಆ
ಆಕಳಿಕೆ
ಆಕಳಿಕೆಯು
ಆಕಸ್ಮಿಕ-ವಾಗಿ
ಆಕಾರ-ದಲ್ಲಿ
ಆಕಾರದ್ದಾಗಿ-ರುತ್ತದೆ
ಆಕಾರ-ವನ್ನು
ಆಕಾಶಕ-ಪೋತಾಸನ
ಆಕಾಶದ
ಆಕಾಶ-ವನ್ನು
ಆಗ
ಆಚರಣೆ-ಗಳು
ಆಡಿಸ-ಬೇಕು
ಆಡಿ-ಸುವು-ದ-ರಿಂದ
ಆಡಿ-ಸು-ವುದು
ಆತಿಥ್ಯ
ಆತ್ಮ
ಆತ್ಮದ
ಆತ್ಮವು
ಆದರೆ
ಆದ್ದ-ರಿಂದ
ಆದ್ದ-ರಿಂ-ದಲೇ
ಆಧಾರ
ಆಧಾರ-ವಿಲ್ಲದೆ
ಆಧಾರ-ಶುದ್ಧಿ
ಆಧಾರ-ಶುದ್ಧಿಯ
ಆನಂದ-ಮಯ-ನಾಗುತ್ತಾನೆ
ಆನಂದ-ವನ್ನು
ಆನೆಯ
ಆನೆ-ಯಂತೆ
ಆಮೆ
ಆಮೆಯ
ಆಯಸ್ಸನ್ನು
ಆಯಾಸ
ಆಯಾಸ-ವನ್ನು
ಆರಂಭದ
ಆರಂಭ-ದಲ್ಲಿ
ಆರಂಭ-ವಾಗುವ
ಆರಂಭಿಸಿ
ಆರಾಧನೆ-ಯನ್ನು
ಆರು
ಆಲಿಂಗನ
ಆಲಿಂಗನ-ದಿಂದ
ಆಲಿಂಗ-ನವು
ಆಲಿಂಗ-ನವೂ
ಆಲಿಂಗ-ನಾಸನ
ಆಲಿಂಗಿಸ-ಬಾರದು
ಆಲಿಂಗಿ-ಸ-ಬೇಕು
ಆಲಿಂಗಿಸಿ
ಆಲಿಂಗಿಸಿ-ದಾಗ
ಆಲಿಂಗಿ-ಸುವ
ಆವರ್ತನೆ-ಗಳನ್ನು
ಆಸಕ್ತಿ
ಆಸನ-ಗಳ
ಆಸನ-ಗಳನ್ನು
ಆಸನ-ಗಳು
ಆಸನವು
ಆಸೆ
ಆಸೆ-ಗಳಿಂದ
ಆಸೆ-ಯಿರುವ-ವಳನ್ನು
ಆಸೆ-ಯಿಲ್ಲದ
ಆಸ್ತಿಕತೆ
ಆಸ್ತಿ-ಗಳೆಲ್ಲವೂ
ಆಹಾರ
ಆಹಾರ-ಗಳೆಂದು
ಆಹಾರದ
ಆಹಾರ-ವನ್ನು
ಆಹಾರ-ವನ್ನೇ
ಇಂಗಳಿಕ
ಇಂತಹ
ಇಂದ್ರನ
ಇಂದ್ರಿಯ
ಇಂದ್ರಿಯ-ಗಳು
ಇಟ್ಟು
ಇಟ್ಟು-ಕೊಂಡು
ಇಟ್ಟು-ಕೊಳ್ಳ-ಬೇಕು
ಇಡ-ಬೇಕು
ಇಡೀ
ಇತರ
ಇತರರ
ಇತರ-ರಿಗೆ
ಇತ್ಯಾದಿ
ಇತ್ಯಾದಿ-ಗಳ
ಇತ್ಯಾದಿ-ಗಳನ್ನು
ಇದಕ್ಕೆ
ಇದನ್ನು
ಇದನ್ನೇ
ಇದರ
ಇದ-ರಿಂದ
ಇದಲ್ಲದೆ
ಇದಾ-ಗಿದೆ
ಇದು
ಇದೆ
ಇದೇ
ಇದ್ದು
ಇನ್ನು
ಇನ್ನೂ
ಇನ್ನೊಂದು
ಇನ್ನೊಬ್ಬ-ಳನ್ನು
ಇಪ್ಪತ್ತ-ನಾಲ್ಕು
ಇಪ್ಪತ್ತು
ಇಪ್ಪತ್ನಾಲ್ಕು
ಇರ-ಬಹುದು
ಇರ-ಬೇಕು
ಇರ-ಲಿಲ್ಲ
ಇರಿ-ಸ-ಬೇಕು
ಇರಿಸಿ
ಇರಿಸಿ-ಕೊಂಡು
ಇರಿಸಿ-ಕೊಳ್ಳ-ಬೇಕು
ಇರುತ್ತವೆ
ಇರುವ
ಇರುವುದ-ರಿಂದ
ಇರು-ವುದಿಲ್ಲ
ಇರು-ವುದೇ
ಇಲ್ಲದೆ
ಇಲ್ಲದೆಯೇ
ಇಲ್ಲಿ
ಇಲ್ಲಿಗೆ
ಇಲ್ಲಿದೆ
ಇಳಿ-ಬಿಟ್ಟಿ-ರುವ
ಇವು
ಇವು-ಗಳ
ಇವು-ಗಳನ್ನು
ಇವು-ಗಳು
ಈ
ಈಗ
ಈಟಿಯ
ಉಂಟಾಗ-ದಿದ್ದರೆ
ಉಂಟಾಗ-ದಿದ್ದಲ್ಲಿ
ಉಂಟಾಗ-ದಿದ್ದಾಗ
ಉಂಟಾಗ-ಬಹು-ದಾದ
ಉಂಟಾಗ-ಬಹುದು
ಉಂಟಾಗುತ್ತದೆ
ಉಂಟಾ-ಗುವ
ಉಂಟಾಗುವ-ವರೆಗೆ
ಉಂಟಾಗು-ವುದಿಲ್ಲ
ಉಂಟಾದ
ಉಂಟಾ-ದರೆ
ಉಂಟಾ-ದಾಗ
ಉಂಟು-ಮಾಡು-ತ್ತದೆ
ಉಗುರಿ-ನಿಂದ
ಉಗುರು-ಗಳ
ಉಗುರು-ಗಳಿಂದ
ಉಚ್ಚೀರ್ಶ-ಕಾಸನ
ಉಜ್ಜ-ಬೇಕು
ಉಜ್ಜಯೀ
ಉಜ್ಜಾಯಿ
ಉಜ್ಜಿದ
ಉಜ್ಜುವ
ಉಜ್ಜು-ವು-ದನ್ನು
ಉಜ್ಜು-ವುದು
ಉಡ್ಯಾಣ
ಉತ್ಕಟಾ-ಸನ
ಉತ್ಕಟಾ-ಸನ-ದಲ್ಲಿ
ಉತ್ತಮ
ಉತ್ತಮ-ವೆಂದು
ಉತ್ತಾನ-ಕೂರ್ಮಾ-ಸನ
ಉತ್ತಾನಾ-ಸನ
ಉತ್ತೇಜನ
ಉತ್ಥಾನೋತ್ತಾ-ನಾಸನ
ಉತ್ಪಾದಕ
ಉತ್ಪೀಡಾ-ಸನ
ಉತ್ಸಾಹವೇ
ಉದರ
ಉದಾಸೀ-ನತೆ
ಉದ್ದ
ಉದ್ದದ
ಉದ್ದ-ವಿರ-ಬೇಕು
ಉದ್ದವು
ಉದ್ದೇಶಕ್ಕಾಗಿ
ಉದ್ದೇಶ-ವಿರ-ಬೇಕು
ಉಪ-ವಾಸ
ಉಪ-ಸಂಹಾರ
ಉಪ್ಪನ್ನು
ಉಪ್ಪು
ಉಬ್ಬಿ-ಕೊಳ್ಳುತ್ತದೆ
ಉಬ್ಬಿ-ಸ-ಬೇಕು
ಉಯ್ಯಾಲೆ
ಉಯ್ಯಾಲೆ-ಯಂತೆ
ಉರಿ
ಉರಿ-ಯನ್ನು
ಉರುಳ-ಬೇಕು
ಉರುಳಾಡ-ಬೇಕು
ಉರುಳುವ
ಉಳಿದ
ಉಳಿ-ಯಂತೆ
ಉಳಿಸಿ-ಕೊಳ್ಳುವ
ಉಳುಮೆ
ಉಷ್ಟ್ರಾಸನ
ಉಷ್ಟ್ರಾಸನ-ದಲ್ಲಿದ್ದಾಗ
ಉಷ್ಣತೆ
ಉಸಿ-ರನ್ನು
ಉಸಿ-ರಾಟ
ಉಸಿ-ರಾಟದ
ಉಸಿರು
ಊದ-ಬೇಕು
ಊರ-ಬೇಕು
ಊರಿ
ಊರ್ಣ-ನಾಭ್ಯಾ-ಸನ
ಊರ್ಧ್ವ-ಪಶ್ಚಿಮತಾ-ನಾಸನ
ಊರ್ಧ್ವಬಿಂದು
ಊರ್ಧ್ವಬಿಂದು-ವಿನ
ಊರ್ಧ್ವ-ಮುಖ-ವಾಗಿ
ಊರ್ಧ್ವ-ರೇತ
ಋಕ್ಷಾ-ಸನ
ಋಷಿ-ಗಳು
ಎಂಟು
ಎಂದರೆ
ಎಂದರೇ
ಎಂದರ್ಥ
ಎಂದಿಗೂ
ಎಂದು
ಎಂದೇ
ಎಂಬ
ಎಂಬು-ದನ್ನರಿತು
ಎಂಬು-ದ-ರಲ್ಲಿ
ಎಂಬುದು
ಎಕ್ಕದ
ಎಚ್ಚರಿಕೆ-ಯಿಂದ
ಎಡ
ಎಡಕ್ಕೆ
ಎಡ-ಗಾಲನ್ನು
ಎಡ-ಗಾಲಿನ
ಎಡಗೈ-ಯನ್ನು
ಎಡಗೈ-ಯಿಂದ
ಎಡ-ದಿಂದ
ಎಣಿ-ಸ-ಬೇಕು
ಎಣಿ-ಸು-ವಾಗ
ಎತ್ತ-ಬೇಕು
ಎತ್ತರ
ಎತ್ತರಕ್ಕೆ
ಎದುರಾ-ಗುತ್ತವೆ
ಎದುರಿನ
ಎದೆಗೆ
ಎದೆಯ
ಎದೆ-ಯನ್ನು
ಎದೆಯು
ಎದ್ದು
ಎನ್ನ-ಲಾಗು-ತ್ತದೆ
ಎರಡನೇ
ಎರಡು
ಎರಡೂ
ಎಲೆ
ಎಲೆ-ಗಳ
ಎಲೆ-ಗಳನ್ನು
ಎಲೆ-ಗಳು
ಎಲೆ-ಯನ್ನು
ಎಲೆ-ಯಿಂದ
ಎಲ್ಲಕ್ಕಿಂತ
ಎಲ್ಲದ-ರಲ್ಲೂ
ಎಲ್ಲವೂ
ಎಲ್ಲಾ
ಎಳಾಯ-ಬೇಕು
ಎಳೆದ
ಎಳೆದು
ಎಳೆದು-ಕೊಂಡಾಗ
ಎಳೆದು-ಕೊಂಡು
ಎಳೆದು-ಕೊಳ್ಳ-ಬೇಕು
ಎಳೆದು-ಕೊಳ್ಳಲು
ಎಳೆದು-ಕೊಳ್ಳುತ್ತಾ
ಎಳೆದು-ಕೊಳ್ಳುತ್ತಾನೆ
ಎಳೆದು-ಕೊಳ್ಳುವ
ಎಳೆದು-ಕೊಳ್ಳು-ವಾಗ
ಎಳೆದು-ಕೊಳ್ಳು-ವುದು
ಎಳೆಯ-ಬೇಕು
ಎಳೆಯು-ತ್ತಾರೋ
ಎಳೆಯುವ
ಎಳೆಯುವು-ದ-ರಿಂದ
ಎಳೆಯು-ವುದು
ಎಳ್ಳೆಣ್ಣೆ
ಎಷ್ಟೆಷ್ಟು
ಎಷ್ಟೇ
ಏಕಾಗ್ರ
ಏರ-ಬೇಕು
ಏರಿ-ಸ-ಬೇಕು
ಏರುವ
ಏಲಕ್ಕಿ-ಯಿಂದ
ಏಳ-ಬೇಕು
ಏಳುತ್ತಿ-ರ-ಬೇಕು
ಏಳು-ವುದು-ಹೀಗೆ
ಒಂಟೆಯ
ಒಂದರ-ಮೇಲೊಂದಿಟ್ಟು
ಒಂದು
ಒಂದೇ
ಒಂದೊಂದೇ
ಒಟ್ಟು
ಒತ್ತಡದ
ಒತ್ತ-ಬೇಕು
ಒತ್ತಿ
ಒದ್ದೆ-ಯಾಗಿ-ರ-ಬೇಕು
ಒನಕೆಯ
ಒಬ್ಬ
ಒಬ್ಬನು
ಒಬ್ಬರ
ಒಮ್ಮೆ
ಒಯ್ದು
ಒಯ್ಯ-ಬೇಕು
ಒಯ್ಯುವುದೇ
ಒಳಕ್ಕೆ
ಒಳ-ಗಿನ
ಒಳಗಿ-ನಿಂದ
ಒಳಗೆ
ಒಳ-ತೊಡೆಯ
ಒಳ-ಸೇರಿಸಿ-ದಾಗ
ಓಂ
ಓಜಸ್ಸು
ಓಡಾಡ-ಬೇಕು
ಓತಿಯ
ಓತಿ-ಯಂತಹ
ಓರೆ-ಯಾದ
ಔಷಧಿ-ಯಂತೆ
ಕಂಕುಳಿನ
ಕಂಕುಳಿ-ನಲ್ಲಿ
ಕಂಥಡೀ
ಕಂದೂ-ಕಾ-ಸನ
ಕಂಬ-ಗ-ಳಂತೆ
ಕಂಬಳಿ
ಕಂಬಿ-ಯಂತಹ
ಕಂಬು-ಕ್ರಿಯೆ-ಯನ್ನು
ಕಚ್ಚಿ
ಕಟ್ಟಿ-ಕೊಂಡು
ಕಟ್ಟಿಗೆ-ಯಿಂದ
ಕಟ್ಟಿ-ರುವ
ಕಟ್ಟು-ಗಳು
ಕಟ್ಟು-ನಿಟ್ಟಿನ
ಕಠಿಣ
ಕಠಿಣ-ವಾದದ್ದಾದ್ದ-ರಿಂದ
ಕಡೆ
ಕಡೆಗೆ
ಕಣ್ಣಿನ
ಕಣ್ಣೀರು
ಕಣ್ಣು-ಗಳನ್ನು
ಕಣ್ಣು-ಗಳಿಗೆ
ಕತ್ತರಿ
ಕತ್ತರಿಯ
ಕತ್ತರಿ-ಸ-ಬೇಕು
ಕತ್ತರಿ-ಸಲು
ಕತ್ತರಿ-ಸು-ವಾಗ
ಕತ್ತರಿ-ಸುವು-ದನ್ನು
ಕತ್ತರಿ-ಸು-ವುದು
ಕತ್ತಿಯ
ಕತ್ತೆಕಿರುಬನ
ಕಪಾಲ
ಕಪಾಲ-ಭಾತಿ
ಕಪಾಲ-ಭಾತಿ-ಯನ್ನು
ಕಪಾಲಾ-ಸನ
ಕಪೋತ-ಕಂಠ
ಕಪೋತ-ನಾಟಕಾ-ಸನ
ಕಪ್ಪು
ಕಪ್ಪೆಯ
ಕಬ್ಬಿಣದ
ಕಬ್ಬಿಣ-ದಿಂದ
ಕಮಠ-ಪೀಠಾಸನ
ಕಮ್ಮಾರನ
ಕರಡಿಯ
ಕರಿ-ಮೆಣಸಿನ
ಕರಿ-ಮೆಣಸು
ಕರುಣೆ
ಕರುಣೆಯ
ಕರೆಯ-ಲಾಗುತ್ತದೆ
ಕರೆಯ-ಲ್ಪಡುವ
ಕರೆಯು-ತ್ತಾರೆ
ಕರೆಯು-ವಾಗ
ಕರೆಯು-ವುದು
ಕರ್ತರಿ
ಕರ್ಪೂರ
ಕರ್ಷ
ಕಲಶ
ಕಲೆಯು
ಕಲ್ಮಶ-ಗಳು
ಕಲ್ಮಶವು
ಕಲ್ಲಿನ
ಕಲ್ಲು
ಕಲ್ಲುಪ್ಪು
ಕಳೆದು-ಕೊಳ್ಳ-ಬಾರದು
ಕಳ್ಳತನ
ಕಳ್ಳರು
ಕಳ್ಳಿಯ
ಕವಿತೆ-ಗಳನ್ನು
ಕಷಾಯ-ಗಳ
ಕಷಾಯ-ವನ್ನು
ಕಸ್ತೂರಿ
ಕಾಂಚ್ಯಾ-ಸನ
ಕಾಂಡ-ಗಳನ್ನು
ಕಾಂಡ-ನಳಿಕೆ-ಗಳ
ಕಾಂಡ-ವನ್ನು
ಕಾಕಾ-ಸನ
ಕಾಕಾ-ಸನ-ದಲ್ಲಿದ್ದಾಗ
ಕಾಕಿ
ಕಾಗೆಯ
ಕಾಚನ್ನು
ಕಾಡಿ-ದರೆ
ಕಾಣಿಸಿ-ಕೊಳ್ಳ-ಬಹುದು
ಕಾಣುವ
ಕಾಣು-ವುದೇ
ಕಾದಂಬಾ-ಸನ
ಕಾಪಾಡಿ-ಕೊಳ್ಳ-ಬೇಕು
ಕಾಪಾಲಿ-ಕರು
ಕಾಮ-ಕ್ರೀಡೆಯ
ಕಾಮ-ಕ್ರೀಡೆ-ಯಲ್ಲಿ
ಕಾಮ-ಕ್ರೀಡೆ-ಯಲ್ಲಿದ್ದರೂ
ಕಾಮದ
ಕಾಮ-ಪ್ರಚೋ-ದಕ
ಕಾಮ-ವನ್ನು
ಕಾಮವು
ಕಾಮೋತ್ಸುಕ-ಳಾದ
ಕಾಮೋ-ದ್ರೇಕ
ಕಾಮೋ-ದ್ರೇಕವು
ಕಾಯಾ
ಕಾಯಿಲೆ-ಗಳನ್ನು
ಕಾಯಿಲೆ-ಗಳಿಂದ
ಕಾಯಿಲೆ-ಗಳಿಲ್ಲ-ದವ-ರಾಗಿ-ರ-ಬೇಕು
ಕಾಯಿಲೆ-ಗಳು
ಕಾರಣ
ಕಾರ್ಯ-ಗಳನ್ನು
ಕಾರ್ಯ-ಗಳು
ಕಾಲ
ಕಾಲ-ದಲ್ಲಿ
ಕಾಲನ್ನು
ಕಾಲಿನ
ಕಾಲಿ-ನಿಂದ
ಕಾಲಿ-ನಿಂದಲೂ
ಕಾಲು
ಕಾಲು-ಗಳ
ಕಾಲು-ಗಳನ್ನು
ಕಾಲ್ತುದಿ-ಯನ್ನು
ಕಾಲ್ಬೆರಳು-ಗಳ
ಕಾಲ್ಬೆರಳು-ಗಳನ್ನು
ಕಾಲ್ಬೆರಳು-ಗಳಿಂದ
ಕಾಳಿನ
ಕಿರು-ಬೆರಳು
ಕಿವಿ
ಕಿವಿ-ಗಳ
ಕಿವಿ-ಗಳನ್ನು
ಕಿವಿ-ಗಳಿಗೆ
ಕಿವಿಯ
ಕಿವಿಯ-ವರೆಗೆ
ಕುಂಟು
ಕುಂಡಲಿನಿ-ಗಳನ್ನು
ಕುಂಡಲಿನಿಯು
ಕುಂಡಲಿನೀ
ಕುಂಬಾರನ
ಕುಂಭಕ
ಕುಂಭಕ-ಗಳ
ಕುಂಭಕ-ಗಳನ್ನು
ಕುಂಭಕ-ಗಳಲ್ಲಿ
ಕುಂಭಕ-ಗಳು
ಕುಂಭಕದ
ಕುಂಭಕ-ವನ್ನು
ಕುಂಭಕಾಃ
ಕುಕ್ಕುಟಾ-ಸನ
ಕುಕ್ಕುಟಾ-ಸನ-ದಲ್ಲಿ-ದ್ದಾಗ
ಕುಕ್ಕುಟಾಸನ-ದಲ್ಲಿದ್ದು
ಕುಕ್ಕುಟೋಡ್ಡಯ-ನಾಸನ
ಕುಗ್ಗಿಸು-ವಂತೆ
ಕುಟೀರದ
ಕುಟೀರ-ದಲ್ಲಿ
ಕುಟೀರ-ವನ್ನು
ಕುಟೀರವು
ಕುಡಿದ
ಕುಡಿದು
ಕುಡಿಯ-ಬೇಕು
ಕುತ್ತಿಗೆ
ಕುತ್ತಿಗೆಯ
ಕುತ್ತಿಗೆ-ಯನ್ನು
ಕುದುರೆ-ಯಂತೆ
ಕುದುರೆಯು
ಕುಬ್ಜಾ-ಸನ
ಕುರಂಟಕ
ಕುರಂಟಕನ
ಕುರಂಟಕರು
ಕುರಿತು
ಕುರಿಯ
ಕುಲಾಲ
ಕುಳಿತಿರು-ವಾಗ
ಕುಳಿತು
ಕುಳಿತು-ಕೊಳ್ಳ-ಬಾರದು
ಕುಳಿತು-ಕೊಳ್ಳ-ಬೇಕು
ಕುಳಿತು-ಕೊಳ್ಳು-ವುದು
ಕುಹರ
ಕುಹರದ
ಕುಹರ-ದ-ವರೆಗೆ
ಕುಹರ-ದೊಳಗೆ
ಕೂಡ
ಕೂಡಿದ
ಕೂಡಿಸಿ
ಕೂದಲಿ-ನಿಂದ
ಕೃಚ್ಛ್ರ
ಕೃಪೆ
ಕೃಷ್ಣನ-ವರೆಗೆ
ಕೆಂಪು
ಕೆಚ್ಚ-ಲನ್ನು
ಕೆನ್ನೆ-ಗಳನ್ನು
ಕೆನ್ನೆ-ಯನ್ನು
ಕೆನ್ನೆ-ಯಿಂದ
ಕೆಮ್ಮಣಿ-ನಿಂದ
ಕೆರೆ-ದಂತೆ
ಕೆರೆ-ಯುವ
ಕೆಲವೇ
ಕೆಲಸ
ಕೆಳ
ಕೆಳಕ್ಕೆ
ಕೆಳ-ಗಿನ
ಕೆಳ-ಗಿನಂತೆ
ಕೆಳ-ಗಿ-ರುವ
ಕೆಳಗೆ
ಕೆಳ-ಭಾಗ
ಕೆಳ-ಭಾಗಕ್ಕೆ
ಕೆಳ-ಭಾಗ-ದಲ್ಲಿ
ಕೆಳ-ಭಾಗದ-ವರೆಗೆ
ಕೆಳ-ಭಾಗ-ವನ್ನು
ಕೆಳ-ಮುಖ-ವಾಗಿ
ಕೇಂದ್ರದ
ಕೇವಲ
ಕೇವಲಃ
ಕೇಶ
ಕೇಶವ
ಕೇಶವ-ನಿಂದ
ಕೇಸರಿ
ಕೈ
ಕೈಗಳ
ಕೈ-ಗಳನ್ನು
ಕೈ-ಗಳಿಂದ
ಕೈಪಿಡಿ
ಕೈಪಿಡಿಯ
ಕೈಪಿಡಿ-ಯನ್ನು
ಕೈ-ಬೆರಳು-ಗಳನ್ನು
ಕೈಯ
ಕೈಯಗಲ
ಕೈ-ಯನ್ನು
ಕೈ-ಯಿಂದ
ಕೊಂಬು
ಕೊಕ್ಕರೆಯ
ಕೊಡಲಿಯ
ಕೊತ್ತಂಬರಿ
ಕೊನೆ-ಗೊಳ್ಳುವ
ಕೊನೆ-ಯಲ್ಲಿ
ಕೊಳಲಿ-ನಂತಹ
ಕೊಳಲಿ-ನಂತೆ
ಕೊಳವೆ
ಕೋರಂಟಕ
ಕೋರಂಟಕನ
ಕೋಲಿನ
ಕೋಲಿ-ನಂತೆ
ಕೌಪೀ-ನಾಸನ
ಕ್ರಮ
ಕ್ರಮ-ದಲ್ಲಿ
ಕ್ರಮ-ದಿಂದ
ಕ್ರಮ-ಬದ್ಧ-ವಾಗಿ
ಕ್ರಮ-ವನ್ನು
ಕ್ರಮೇಣ
ಕ್ರಿಮಿ
ಕ್ರಿಯೆ
ಕ್ರಿಯೆ-ಗಳ
ಕ್ರಿಯೆ-ಗಳನ್ನು
ಕ್ರಿಯೆ-ಗಳಿಂದ
ಕ್ರಿಯೆಗೆ
ಕ್ರಿಯೆಯ
ಕ್ರಿಯೆ-ಯನ್ನು
ಕ್ರಿಯೆ-ಯಲ್ಲಿ
ಕ್ರಿಯೆ-ಯಿಂದ
ಕ್ರಿಯೆಯು
ಕ್ರೀಡೆ
ಕ್ರೀಡೆ-ಗಳಲ್ಲಿ
ಕ್ರೀಡೆ-ಗಳೂ
ಕ್ರೌಂಚಾ-ಸನ
ಕ್ಷಣ-ದಲ್ಲಿ
ಕ್ಷಣಾರ್ಧ-ದಲ್ಲಿ
ಕ್ಷೀಣಿ-ಸಿ-ದರೂ
ಖಂಡಿತ-ವಾಗಿಯೂ
ಖಚಿತ-ವಾಗಿ
ಖಡ್ಗ-ಮೃಗ
ಖಡ್ಗ-ಮೃಗದ
ಖಡ್ಗಾ-ಸನ
ಖಾಲಿ
ಖೇಚರಿ
ಗಂಟನ್ನು
ಗಂಟಲನ್ನು
ಗಂಟಲಿನ
ಗಂಟಲಿ-ನಂತೆ
ಗಂಟಲಿನ-ವರೆಗೆ
ಗಂಟಲು
ಗಂಟೆ-ಗಳ
ಗಂಟೆ-ಗಳ-ವರೆಗೆ
ಗಂಭೀರ
ಗಜ-ಕರಣಿ
ಗಜ-ಸಾಧನಾ
ಗಜಾ-ಸನ
ಗಜಾ-ಸನ-ದಲ್ಲಿದ್ದಾಗ
ಗಜಾ-ಸನ-ದಲ್ಲಿ-ರುವಂತೆಯೇ
ಗಟ್ಟಿ-ಯಾಗಿ
ಗಟ್ಟಿ-ಯಾದ
ಗಣೇಶ
ಗಣೇಶನ
ಗಣೇಶ-ನಿಗೆ
ಗರಿಷ್ಠ
ಗರುಡಾ-ಸನ
ಗಲ್ಲ-ವನ್ನು
ಗಳಿಸು-ತ್ತಾನೆ
ಗಹನ
ಗಾತ್ರ-ದಷ್ಟಿ-ರ-ಬೇಕು
ಗಾಯಕಿ-ಯರಾಗಿ-ರ-ಬೇಕು
ಗಾಳಿ-ಯನ್ನು
ಗಾಳಿ-ಯಲ್ಲಿ
ಗಾಳಿಯು
ಗಿಡ
ಗಿಡದ
ಗಿಡ-ಮೂಲಿಕೆ-ಗಳ
ಗಿಡು-ಗದ
ಗಿರ-ಗಿರನೆ
ಗಿರಣಿ
ಗಿಳಿಯ
ಗುಣ-ಗಳ
ಗುಣ-ಗಾನ-ವನ್ನು
ಗುಣ-ಪಡಿಸಿ-ಕೊಳ್ಳಲು
ಗುದಕಂಬು-ಕ್ರಿಯಾ
ಗುದ-ದ್ವಾರ
ಗುದ-ದ್ವಾರದ
ಗುದ-ದ್ವಾರ-ದಲ್ಲಿ
ಗುದ-ದ್ವಾರ-ದೊಳಗೆ
ಗುದ-ದ್ವಾರ-ವನ್ನು
ಗುದ-ನಾಳ-ದಲ್ಲಿ
ಗುಪ್ತ
ಗುಬ್ಬಚ್ಚಿಯ
ಗುರು-ಗಳಿಗೆ
ಗುರು-ಗಳು
ಗುರುತಿ-ಸುವ
ಗುರುತು
ಗುರುಭ್ಯೋ
ಗುರು-ವಿನ
ಗುರುವಿ-ನಿಂದ
ಗುರುವಿ-ನೆಡೆ-ಗಿನ
ಗುರು-ಸೇವೆ
ಗುಳಿ-ಗಳನ್ನು
ಗುಳಿ-ಯಲ್ಲಿ
ಗುಳ್ಳೆ-ಗಳ-ಲ್ಲಿನ
ಗುಳ್ಳೆ-ಗಳು
ಗೂನು
ಗೂಳಿ
ಗೃಹಸ್ಥರು
ಗೇಣಿ-ಗಿಂತ
ಗೇಣು
ಗೇಣು-ಗಳಷ್ಟು
ಗೋಡೆ
ಗೋಡೆಗೆ
ಗೋಡೆ-ಯನ್ನು
ಗೋಡೆ-ಯಿಂದ
ಗೋಧಿ
ಗೋರಕ್ಷ-ನಾಥರು
ಗೌಜ-ಲಕ್ಕಿಯ
ಗೌರ-ವರ್ಣದ-ವರು
ಗ್ರಂಥ-ದಲ್ಲಿ
ಗ್ರಾಮ್
ಗ್ರಾಹಾ-ಸನ
ಘಟಕ
ಘಟಕ-ಗಳಿಂದ
ಘಟಕ-ವಿಲ್ಲದೆಯೇ
ಘರ್ಷಣೆ-ಯಿಂದ
ಚಂದ್ರ
ಚಂದ್ರನ
ಚಂದ್ರ-ನನ್ನು
ಚಂದ್ರ-ರನ್ನು
ಚಕ್ಕೆ-ಯನ್ನು
ಚಕ್ರ-ಗಳ
ಚಕ್ರದ
ಚಕ್ರ-ದಿಂದ
ಚಕ್ರಾ-ಸನ
ಚಟಕಾ-ಸನ
ಚಟಕಾ-ಸನ-ದಂತೆಯೇ
ಚಟುವಟಿಕೆ-ಯಿಂದ
ಚರ್ಮ-ವನ್ನು
ಚಲನ
ಚಲ-ನವು
ಚಲನೆ
ಚಲಿಸ-ಬೇಕು
ಚಲಿಸಿ-ದರೆ
ಚಲಿ-ಸುವ
ಚಲಿಸು-ವಂತೆ
ಚಲಿಸು-ವುದು
ಚಶಕ
ಚಾಂದ್ರಾಯಣ-ದಂತಹ
ಚಾಚ-ಬೇಕು
ಚಾಚಿ
ಚಾಚಿದ
ಚಾಚುತ್ತಾ
ಚಾರ್ಯ
ಚಾಲನ
ಚಾಲನೆ-ಯನ್ನು
ಚಾಷ್ಟ
ಚಿಂತೆ-ಯಿಲ್ಲದೇ
ಚಿಗುರಿ-ನಷ್ಟು
ಚಿಟಿಕೆ
ಚಿತ್ರ-ಮೂಲದ
ಚಿನ್ನ
ಚಿನ್ನ-ದಿಂದ
ಚಿನ್ಮುದ್ರೆ-ಯಲ್ಲಿ-ರಿಸಿ
ಚಿಪ್ಪಿನ
ಚುಂಬಿ-ಸು-ವುದು
ಚೂರನ್ನು
ಚೆಂಡಿನ
ಚೆನ್ನಾಗಿ
ಚೌರಾಸಿ
ಛತ್ರಾ-ಸನ
ಛತ್ರಿಯ
ಜಗತ್ತಿನ
ಜನನ-ಮರಣ-ಗಳ
ಜನನಾಂಗ
ಜನನಾಂಗದ
ಜನ್ಮ-ಗಳಲ್ಲಿ
ಜಪ
ಜಪದ
ಜಪಿ-ಸ-ಬೇಕು
ಜಪಿ-ಸುವುದೇ
ಜಯಿ-ಸಲು
ಜಾಗದ
ಜಾಗ-ದಿಂದ
ಜಾಗ-ರೂಕತೆ-ಯಿಂದ
ಜಾಗೃತ-ಗೊಳಿಸಿ
ಜಾಗೃತ-ಗೊಳಿ-ಸುತ್ತದೆ
ಜಾರುವ
ಜಾಲಂಧರ
ಜಿಂಕೆಯ
ಜಿಗಿ-ಯುತ್ತಾ
ಜಿರಳೆಯ
ಜಿಹ್ವಾ-ಮೈಥುನ
ಜಿಹ್ವಾ-ಮೈಥುನ-ವನ್ನೂ
ಜೀರ್ಣ-ವಾಗುತ್ತದೆ
ಜೀವಿ-ಗಳ
ಜುಟ್ಟಿನ-ವರೆಗೆ
ಜೇಂಕಾರ-ದಂತಹ
ಜೇಡದ
ಜೋಡಣೆ-ಯನ್ನು
ಜೋಡಿ-ಸ-ಬೇಕು
ಜೋಡಿಸಿ
ಜೋಡಿಸಿ-ರ-ಬೇಕು
ಜ್ಞಾನ-ವಿಲ್ಲದೆ
ಜ್ಞಾನವು
ಜ್ಞಾನಿ-ಯಾದ
ಜ್ಞಾನಿ-ಯಾದ-ವನು
ಜ್ವರ
ಜ್ವರ-ವನ್ನು
ಜ್ವರ-ವುಂಟಾ-ಗುತ್ತದೆ
ತಂತೀ-ವಾದ್ಯ-ಗಳನ್ನು
ತಂತ್ರ
ತಂತ್ರ-ಗಳಿಂದ
ತಂತ್ರ-ಗಳು
ತಂತ್ರ-ದಂತಹ
ತಂತ್ರ-ವನ್ನು
ತಂತ್ರವು
ತಂತ್ರ-ವೆಂದು
ತಂದರು
ತಂದು
ತಕ್ಕಂತೆ
ತಕ್ಷಣ
ತಕ್ಷಣವೇ
ತಗಲಿ-ರಲಿ
ತಗಲು-ವಂತಿರ-ಬೇಕು
ತಗಲು-ವಂತೆ
ತಡೆದು-ಕೊಳ್ಳು-ವಷ್ಟು
ತಡೆ-ಹಿಡಿದ
ತಡೆ-ಹಿಡಿದು
ತಡೆ-ಹಿಡಿಯ-ಬೇಕು
ತಡೆ-ಹಿಡಿಯು-ವುದು
ತತ್ವ-ಗಳು
ತಥಾ
ತದ-ನಂತರ
ತನ್ನ
ತನ್ನನ್ನು
ತನ್ನೆಲ್ಲಾ
ತಪಸ್ವಿ-ಗಳು
ತಪಸ್ಸು
ತಪ್ಪಿ-ಸಲು
ತಬ್ಬಿ
ತಬ್ಬಿ-ಕೊಂಡು
ತಯಾರಿ-ಸ-ಬಹುದು
ತಯಾರಿ-ಸ-ಬೇಕು
ತಯಾರಿ-ಸ-ಲಾಗುತ್ತದೆ
ತಯಾ-ರಿಸಿ
ತಯಾ-ರಿಸಿ-ಕೊಳ್ಳ-ಬೇಕು
ತಯಾರಿ-ಸುವ
ತರಕ್ಷ್ವಾ-ಸನ
ತರ-ಬೇಕು
ತಲುಪ-ದಿದ್ದಾಗ
ತಲುಪಿ-ದಾಗ
ತಲುಪು-ತ್ತದೆ
ತಲೆ
ತಲೆಯ
ತಲೆ-ಯನ್ನು
ತಲೆ-ಯ-ವರೆಗೆ
ತಲೆ-ಯಿಂದ
ತಲೆ-ಸುತ್ತು
ತಳ-ಭಾಗ-ದಲ್ಲಿ
ತಳ್ಳ-ಬೇಕು
ತಾಂಡವಾ-ಸನ
ತಾಗಿ-ಸ-ಬೇಕು
ತಾಗಿ-ಸ-ಬೇಕು-ನಂತರ
ತಾಗಿಸಿ
ತಾಗಿ-ಸುವ
ತಾಜಾ
ತಾತ್ಕಾಲಿಕ
ತಾಳ್ಮೆ
ತಾಳ್ಮೆ-ಮತ್ತು
ತಿಂಗಳ
ತಿಂಗಳಲ್ಲಿ
ತಿಂಗಳೊಳಗೆ
ತಿತ್ತಿರ್ಯಾ-ಸನ
ತಿರುಗಿ-ದಾಗ
ತಿರುಗಿ-ಸ-ಬೇಕು
ತಿರುಗಿ-ಸುವ
ತಿರುಗಿ-ಸುವು-ದಕ್ಕೆ
ತಿರುಗು-ಮುರುಗು
ತಿರ್ಯನ್ನೌಕಾ-ಸನ
ತಿರ್ಯನ್ನೌಕಾ-ಸನ-ದಲ್ಲಿ-ರುವಂತೆಯೇ
ತಿಳಿದಿ-ರು-ವಂತೆ
ತಿಳಿಯ-ಬೇಕು
ತಿಳಿಯಲು
ತಿಳಿಯುವ
ತೀರ್ಥ-ಯಾತ್ರೆ
ತೀವ್ರ
ತೀವ್ರ-ವಾದ
ತುಂಬಿ-ಕೊಂಡು
ತುಂಬಿಸಿ
ತುಂಬು-ವಂತೆ
ತುಟಿ-ಗಳನ್ನು
ತುದಿ-ಗಳನ್ನು
ತುದಿ-ಗಳು
ತುದಿಗೆ
ತುದಿಯ
ತುದಿ-ಯನ್ನು
ತುದಿ-ಯಿಂದ
ತುಪ್ಪ
ತುರಿಕೆ
ತುರಿಕೆಯ
ತುರಿಕೆಯು
ತೂಕದ
ತೂಗುವ
ತೂರಿಸಿ
ತೃಣಜಲೂಕಾ-ಸನ
ತೃಪ್ತಿ
ತೆಗೆದು-ಕೊಂಡು
ತೆಗೆದು-ಕೊಳ್ಳ-ಬೇಕು
ತೆಗೆದು-ಹಾಕ-ಬೇಕು
ತೆಗೆದು-ಹಾಕಲು
ತೆರೆದು
ತೆಳ್ಳ-ಗಾಗುತ್ತದೆ
ತೆಳ್ಳ-ಗಾ-ದರೂ
ತೇಗುವ
ತೇಜಸ್ಸು
ತೇರಿನ
ತೊಂದರೆ-ಯಾ-ದರೂ
ತೊಟ್ಟೂ
ತೊಡಗ-ಬಾರದು
ತೊಡಗ-ಬೇಕು
ತೊಡ-ಗಲು
ತೊಡೆ
ತೊಡೆ-ಗಳ
ತೊಡೆ-ಗಳನ್ನು
ತೊಡೆ-ಗಳು
ತೊಡೆಗೆ
ತೊಡೆಯ
ತೊಳೆದ
ತೋರು-ಬೆರಳು
ತೋರು-ಬೆರಳು-ಗಳನ್ನು
ತೋರು-ಬೆರಳು-ಗಳಿಂದ
ತೋರುವ
ತೋಳದ
ತೋಳಿನ
ತೋಳಿ-ನಿಂದ
ತೋಳು-ಗಳ
ತೋಳು-ಗಳನ್ನು
ತೋಳು-ಗಳು
ತ್ಯಜಿ-ಸ-ಬೇಕು
ತ್ಯಜಿ-ಸುವ
ತ್ರಾಟಕ
ತ್ರಿಕೂಟ
ತ್ರಿಕೋನ
ತ್ರಿವಿಕ್ರ-ಮಾಸನ
ದಂಡ-ದಂತೆ
ದಂಡಾ-ಸನ
ದಂತ
ದಪ್ಪ-ವಾಗಿದ್ದು
ದಪ್ಪ-ವಿ-ರುವ
ದಪ್ಪವು
ದಾಟುವ
ದಾರದ
ದಾರ-ವನ್ನು
ದಾರವು
ದಾರಿ
ದಿಂಬಿನ
ದಿಟ್ಟಿಸಿ
ದಿನಕ್ಕೆ
ದಿನದ-ವರೆಗೆ
ದೀಕ್ಷೆ
ದೀರ್ಘಾಯುಷ್ಯವು
ದುಂಡಗೆ
ದುಂಬಿಯ
ದುಷ್ಟ-ರಿಂದ
ದೂರ-ದಲ್ಲಿ
ದೂರ-ವಿ-ರ-ಬೇಕು
ದೂರಾದ-ವನ
ದೃಢವಾಗುತ್ತದೆ
ದೃಶದಾಸನ
ದೇವತೆ-ಗಳ
ದೇವರ
ದೇಹ
ದೇಹಕ್ಕೆ
ದೇಹದ
ದೇಹಲ್ಯುಲ್ಲಂಘನಾಸನ
ದೇಹ-ವನ್ನು
ದೈವ-ಕೃಪೆ-ಯಿಂದ
ದೈವಾರಾಧನೆ
ದೈವಾರಾಧನೆಯ
ದೈಹಿಕ
ದೋಣಿ
ದೋಣಿಯ
ದೋಷ-ಗಳ
ದೋಷ-ಗಳು
ದ್ರವ-ಗಳನ್ನು
ದ್ರವ-ವನ್ನು
ದ್ರವ್ಯ-ಗಳನ್ನು
ದ್ವಾದ-ಶಾಕ್ಷರಿ
ದ್ವಾರಕ್ಕೆ
ದ್ವಾರದೊಳಗೆ
ದ್ವಾರ-ವನ್ನು
ದ್ವಿತೀಯ
ದ್ವಿಶೀರ್ಷಾಸನ
ಧಕ್ಕೆಯುಂಟು
ಧನುರಾಸನ
ಧನ್ವಂತರಿಯ
ಧರ್ಮ
ಧಾರಣೆ
ಧೈರ್ಯ
ಧೈರ್ಯ-ದಿಂದ
ಧೌತಿ
ಧೌತಿ-ಗಳಂತಹ
ಧ್ಯಾನ-ದಿಂದ
ಧ್ಯಾನಿ-ಸ-ಬೇಕು
ಧ್ರುವಾಸನ
ಧ್ವಜಾಸನ
ಧ್ವನಿಯ
ಧ್ವನಿ-ಯಲ್ಲಿ
ನಂತರ
ನಂಬಿಕೆಯಿ-ರುವುದೇ
ನಕಲಿ
ನಖ
ನಡು-ವಿನ
ನಡುವೆ
ನಡೆ-ಸ-ಬೇಕು
ನಡೆ-ಸಿದ
ನಮಃ
ನಮನ-ಗಳು
ನಮೋ
ನಯ-ವಾಗಿ-ರ-ಬೇಕು
ನಯ-ವಾದ
ನರಕಾಸನ
ನರಕಾಸನ-ದಲ್ಲಿ-ರುವಾ-ಗಲೇ
ನರ-ಗಳ
ನರ-ವನ್ನು
ನಲವತ್ತೆಂಟು
ನಳಿಕೆ
ನಳಿಕೆ-ಗಳನ್ನು
ನಳಿಕೆ-ಗಳೊಂದಿ-ಗಿನ
ನಳಿಕೆಯ
ನಳಿಕೆ-ಯನ್ನು
ನಳಿಕೆಯು
ನಳಿಕೆ-ಯೊಂದಿಗೆ
ನಳಿಕೆ-ಯೊಳಗೆ
ನವಣೆ
ನವಿಲಿನ
ನಶ್ವರತೆ-ಯನ್ನು
ನಷ್ಟವಾಗ-ದಂತೆ
ನಷ್ಟವಾ-ಗಲೇಬಾರದು
ನಾಡಿ-ಗಳ
ನಾಡಿ-ಗಳು
ನಾಡಿ-ಯನ್ನು
ನಾಡೀಮಂಡಲವು
ನಾದ-ವನ್ನು
ನಾನು
ನಾಭಿ
ನಾಭಿಯ
ನಾಭಿ-ಯನ್ನು
ನಾಭಿ-ಯ-ವರೆಗೆ
ನಾಮ
ನಾಮ-ಗಳನ್ನು
ನಾಮಸ್ಮರಣೆ
ನಾಮಸ್ಮರಣೆ-ಯಿಂದ
ನಾಯಿಯ
ನಾರದಾಸನ
ನಾಲಿಗೆಯ
ನಾಲಿಗೆ-ಯನ್ನು
ನಾಲಿಗೆಯು
ನಾಲ್ಕು
ನಾಳ-ವನ್ನು
ನಾಶ-ವಾಗು-ತ್ತವೆ
ನಾಸಿಕಾ
ನಿಂತಿ-ರುವ
ನಿಂತಿ-ರು-ವಾಗ
ನಿಂತು
ನಿಂತು-ಹೋಗುತ್ತ-ದೆಯೋ
ನಿಕಟ
ನಿಜ-ವಾದ
ನಿತ್ಯ
ನಿದ್ರೆಗೆ
ನಿಧಾನ-ವಾಗಿ
ನಿಮಿಷ-ಗಳಲ್ಲಿ
ನಿಮಿಷ-ದಿಂದ
ನಿಯಂತ್ರಣ
ನಿಯಮ
ನಿಯಮ-ಗಳ
ನಿಯಮ-ಗಳನ್ನು
ನಿಯಮ-ಗಳಿ-ರು-ವುದಿಲ್ಲ
ನಿಯಮ-ಗಳು
ನಿರಂತರ
ನಿರಂತರ-ವಾಗಿ
ನಿರರ್ಥಕ-ವೆನಿ-ಸಿ-ದಾಗ
ನಿರ್ಜನ
ನಿರ್ದೇಶಿ-ಸ-ಬೇಕು
ನಿರ್ಮಲ-ವಾದ
ನಿಲ್ಲ-ಬೇಕು
ನಿಲ್ಲಿ-ಸ-ಬೇಕು
ನಿಲ್ಲುವ-ವರೆಗೆ
ನಿವಾರಣೆ-ಯಾಗುತ್ತವೆ
ನಿವಾರಣೆ-ಯಾದ
ನಿವಾರಿ-ಸಲು
ನಿವಾರಿ-ಸಿ-ಕೊಳ್ಳಲು
ನಿಶ್ಚಲ-ವಾಗಿ
ನಿಷ್ಠೆ-ಯಿಂದ
ನೀಡ-ದಿದ್ದಾಗ
ನೀಡ-ಬೇಕು
ನೀಡಲು
ನೀಡು-ತ್ತದೆ
ನೀಡುತ್ತ-ದೆಯೋ
ನೀರನ್ನು
ನೀರಸ-ವೆನಿ-ಸಿ-ದಾಗ
ನೀರಾಸಕ್ತ
ನೀರಿಗೆ
ನೀರಿ-ನಲ್ಲಿ
ನೀರಿ-ನಿಂದ
ನೀರಿ-ನೊಂದಿಗೆ
ನೀರಿ-ರ-ಬೇಕು
ನೀರು
ನುಂಗ-ಬೇಕು
ನುಂಗಿ
ನುಂಗಿದ
ನುಂಗುವ
ನುಂಗು-ವುದು
ನುಡಿಸು-ವುದ-ರಲ್ಲಿ
ನೂರಕ್ಕೂ
ನೂರಾರು
ನೃತ್ಯ
ನೃತ್ಯದ
ನೃತ್ಯಾ-ಸನ-ದಲ್ಲಿದ್ದಾಗ
ನೆನೆಸ-ಬೇಕು
ನೆನೆಸಿಟ್ಟು
ನೆರ-ಳನ್ನು
ನೆರವಿ-ನಿಂದ
ನೆಲಕ್ಕೆ
ನೆಲದ
ನೆಲ-ದಿಂದ
ನೆಲ-ವನ್ನು
ನೆಲೆ-ನಿಲ್ಲು-ವುದಿಲ್ಲ
ನೆಲೆ-ಸುತ್ತದೆ
ನೇಗಿ-ಲಿನ
ನೇಗಿ-ಲಿ-ನಿಂದ
ನೇತಿ
ನೇರ-ವಾಗಿ
ನೇರ-ವಾಗಿ-ರ-ಬೇಕು
ನೇರ-ವಾಗಿ-ರಿಸಿ
ನೇರ-ವಾಗಿ-ಸ-ಬೇಕು
ನೇರ-ವಾಗಿಸಿ
ನೇರ-ವಾದ
ನೇವ-ರಿಸಿ-ಕೊಳ್ಳ-ಬೇಕು
ನೈಜ
ನೈತಿಕ
ನೈಪುಣ್ಯ
ನೋಡ-ಬೇಕು
ನೋಡಲು
ನೋಡು-ವಾಗ
ನೋಡು-ವುದ-ರಿಂದಲೂ
ನೋಡು-ವುದು
ನೋವನ್ನು
ನೋವ-ನ್ನುಂಟು
ನೋವಿ-ನೊಂದಿಗೆ
ನೋವು
ನೌಕಾ-ಸನ
ನೌಕಾ-ಸನ-ದಲ್ಲಿ-ರುವಾ-ಗಲೇ
ನೌಲಿ
ನ್ಯುಬ್ಜಾ-ಸನ
ಪಂಗು-ಕುಕ್ಕುಟಾ-ಸನ
ಪಂಗು-ಮಯೂರಾ-ಸನ
ಪಂಚೇಂದ್ರಿಯ-ಗಳ
ಪಕ್ಕಕ್ಕೆ
ಪಕ್ಕ-ದಲ್ಲಿ
ಪಕ್ಕೆಲು-ಬು-ಗಳ
ಪಕ್ವ-ವಾಗು-ತ್ತವೆ
ಪಚ್ಚ
ಪಟ್ಟಿಯ
ಪಠಣೆ
ಪಠಿ-ಸ-ಬೇಕು
ಪಠಿ-ಸುತ್ತಿ-ರ-ಬೇಕು
ಪಠಿಸು-ವುದ-ರಿಂದ
ಪಠಿಸು-ವುದ-ರಿಂದಲೂ
ಪಡದಿ-ರುವುದೇ
ಪಡು-ವುದೇ
ಪಡೆ-ದ-ವನು
ಪಡೆದು
ಪಡೆಯ-ಬಹುದು
ಪಡೆಯ-ಬೇಕು
ಪಡೆಯ-ಲಾರ
ಪಡೆ-ಯಲು
ಪಡೆಯು-ತ್ತಾನೆ
ಪತನ-ವಾ-ಗುವ
ಪತನ-ವಾ-ದರೆ
ಪತಿ
ಪತ್ನಿ
ಪಥ್ಯದ
ಪದಾರ್ಥ-ಗಳಿಂದ
ಪದೇ
ಪದ್ಧತಿ
ಪದ್ಧತಿಯ
ಪದ್ಮಾಸನ
ಪದ್ಮಾಸನ-ದಲ್ಲಿ
ಪದ್ಮಾಸನ-ದಲ್ಲಿದ್ದಾಗ
ಪದ್ಮಾಸನ-ದಲ್ಲಿ-ರುವ
ಪರಮಾ-ನಂದ-ವಾಗಿ
ಪರರ
ಪರ-ಶ್ವಾಸನ
ಪರಸ್ಪರ
ಪರಿಗಣಿಸ-ಲಾಗು-ತ್ತದೆ
ಪರಿಘಾ-ಸನ
ಪರಿ-ಣಾಮ
ಪರಿಣಿತಿ
ಪರಿಣಿತಿಯ
ಪರಿಣಿತಿ-ಯನ್ನು
ಪರಿಹರಿ-ಸಿ-ಕೊಳ್ಳಲು
ಪರೋಪಕಾರ
ಪರೋಶ್ಣ್ಯಾಸನ
ಪರ್ಪಟಾಸನ
ಪರ್ಯಂಕಾಸನ
ಪರ್ಯಂಕಾಸನ-ದಲ್ಲಿ-ರುವಾ-ಗಲೇ
ಪರ್ವತಾಸನ
ಪಲ
ಪಲ-ಗಳಷ್ಟು
ಪವಿತ್ರಗೊಳಿಸುವುದೇ
ಪಶ್ಚಾತಾಪ
ಪಶ್ಚಿಮತಾನಾಸನ
ಪಶ್ಚಿಮತಾನಾಸನದ
ಪಶ್ಚಿಮತಾನಾಸನ-ವನ್ನು
ಪಾತ್ರರಾಗುವುದೇ
ಪಾದ
ಪಾದ-ಗಳ
ಪಾದ-ಗಳನ್ನಿಟ್ಟು
ಪಾದ-ಗಳನ್ನು
ಪಾದದ
ಪಾದ-ದಿಂದಲೂ
ಪಾದ-ರಕ್ಷೆಯ
ಪಾದ-ರಸ-ವನ್ನು
ಪಾದ-ವನ್ನು
ಪಾದುಕಾಸನ
ಪಾಪ
ಪಾರಿವಾಳದ
ಪಾರ್ಶ್ವದ
ಪಾಲಿ-ಸ-ಬೇಕು
ಪಾಲಿ-ಸುವ
ಪಾಲಿಸುವುದೇ
ಪಾಶಾಸನ
ಪಾಷಂಡಿ-ಗಳಲ್ಲಿ
ಪೀಠ
ಪೀಠಿಕೆ
ಪೀಡಿತರಾದ-ವರು
ಪುಡಿ
ಪುಡಿ-ಯನ್ನು
ಪುಣ್ಯ
ಪುಣ್ಯಕ್ಷೇತ್ರ-ಗಳಿಗೆ
ಪುಣ್ಯದ
ಪುಣ್ಯ-ವನ್ನು
ಪುಣ್ಯವು
ಪುನಃ
ಪುರುಷ-ರನ್ನು
ಪುರುಷರು
ಪೂಜ್ಯ
ಪೂರ್ಣ
ಪೂರ್ಣ-ಗೊಳಿಸ-ಬಹುದು
ಪೂರ್ಣ-ಗೊಳ್ಳುತ್ತದೆ
ಪೂರ್ಣ-ವಾಗಿ
ಪೂರ್ಣ-ವಾಯಿತು
ಪೂರ್ತಿ-ಯಾಗಿ
ಪೂರ್ವ-ಸಿದ್ಧತಾ
ಪೃಷ್ಠಕ್ಕೆ
ಪೃಷ್ಠದ
ಪೃಷ್ಠವನ್ನಿಟ್ಟು
ಪೃಷ್ಠ-ವನ್ನು
ಪೃಷ್ಠವು
ಪೆಟ್ಟು
ಪೆಟ್ಟು-ಗಳನ್ನು
ಪೋಷಣೆಯ
ಪ್ರಕಾರ
ಪ್ರಕ್ರಿಯೆ-ಗಳನ್ನು
ಪ್ರಕ್ರಿಯೆ-ಯಲ್ಲಿ
ಪ್ರಕ್ರಿಯೆಯು
ಪ್ರಕ್ಷಾ-ಲನೆ
ಪ್ರಚೋದನೆ
ಪ್ರಚೋದಿತ-ಳಾದ
ಪ್ರಚೋದಿತ-ವಾ-ದಾಗ
ಪ್ರಚೋದಿ-ಸ-ಬೇಕು
ಪ್ರತಿ-ದಿನ
ಪ್ರತಿ-ಯೊಬ್ಬ-ಳನ್ನು
ಪ್ರದಕ್ಷಿಣಾ-ಕಾರ-ವಾಗಿ
ಪ್ರದೇಶ-ದಲ್ಲಿ
ಪ್ರಭಾವ
ಪ್ರಭಾವ-ದಿಂದ
ಪ್ರಭುತ್ವ
ಪ್ರಮಾಣದ
ಪ್ರಮಾಣ-ದಲ್ಲಿ
ಪ್ರಮುಖ
ಪ್ರಯತ್ನ-ದೊಂದಿಗೆ
ಪ್ರಯತ್ನಿ-ಸುತ್ತವೆ
ಪ್ರಯಾಣಿ-ಸು-ವುದೇ
ಪ್ರವೇಶ
ಪ್ರವೇಶಿ-ಸುತ್ತದೆ
ಪ್ರಶಂಸಿ-ಸುತ್ತೇನೆ
ಪ್ರಸಿದ್ಧ-ವಾ-ಗಿದೆ
ಪ್ರಸಿದ್ಧ-ವಾ-ಗಿವೆ
ಪ್ರಸ್ಥ
ಪ್ರಸ್ಥ-ಗಳಷ್ಟು
ಪ್ರಸ್ಥ-ಗಳು
ಪ್ರಾಕೃತದ
ಪ್ರಾಣ
ಪ್ರಾಣ-ವಾಯುವು
ಪ್ರಾಣವು
ಪ್ರಾಣ-ಶಕ್ತಿ
ಪ್ರಾಣಾಯಾಮ-ಗಳು
ಪ್ರಾಣಾಯಾಮದ
ಪ್ರಾಣಿ-ಗಳ
ಪ್ರಾರಂಭ-ವಾ-ಗುವ
ಪ್ರಾರಂಭಿ-ಸ-ಬೇಕು
ಪ್ರಾರಂಭಿಸಿ
ಪ್ರಾರಂಭಿಸಿ-ದಾಗ
ಪ್ರೀತಿ
ಪ್ರೀತಿಯೇ
ಪ್ರೇಂಖಾ-ಸನ
ಪ್ರೇಂಖಾ-ಸನ-ದಲ್ಲಿದ್ದಾಗ
ಪ್ರೇಮ-ಪೂರ್ವಕ
ಫಲ
ಫಲ-ವನ್ನು
ಫಲ-ವಾಗಿ
ಫೂಟ್
ಬಂದ
ಬಂಧ
ಬಂಧ-ನ-ದಿಂದ
ಬಂಧ-ನ-ವನ್ನು
ಬಂಧ-ನವು
ಬಂಧಿಸ-ಬೇಕು
ಬಕಾ-ಸನ
ಬಗ್ಗೆ
ಬಗ್ಗೆ-ಯಾ-ಗಲಿ
ಬಜೆ
ಬಟ್ಟಲು
ಬಟ್ಟೆಯ
ಬಟ್ಟೆ-ಯನ್ನು
ಬಟ್ಟೆ-ಯಿಂದ
ಬಟ್ಟೆಯು
ಬಣ್ಣದ
ಬದಲಾ-ಗುತ್ತದೆ
ಬದಲಿಗೆ
ಬದಿ-ಗಳ
ಬದಿ-ಯಲ್ಲಿ
ಬದಿ-ಯಲ್ಲೂ
ಬದ್ಧ
ಬಯ-ಸುವ
ಬಯ-ಸುವ-ವನು
ಬರ-ಬಾರದು
ಬರುತ್ತದೆ
ಬರುವ
ಬರುವಂತಾ-ದಾಗ
ಬರು-ವಂತೆ
ಬರು-ವ-ವರೆಗೆ
ಬರು-ವುದೇ
ಬರೆ-ಯಲಾ-ಗಿದೆ
ಬಲ
ಬಲಕ್ಕೆ
ಬಲ-ಗಾಲಿನ
ಬಲಗೈ
ಬಲಗೈ-ಯನ್ನು
ಬಲಗೈ-ಯಿಂದ
ಬಲಗೈ-ಯೊಂದಿಗೆ
ಬಲ-ದಿಂದ
ಬಲ-ಪಡಿ-ಸಲು
ಬಲ-ಪ್ರಯೋಗ
ಬಲ-ವನ್ನು
ಬಲ-ವಾಗಿ
ಬಲವು
ಬಲ-ಶಾಲಿ-ಯಾಗು-ತ್ತಾನೆ
ಬಳಸ-ಬಾರದು
ಬಳಸ-ಬೇಕು
ಬಳಸಿ
ಬಳ-ಸುವ
ಬಳ-ಸುವು-ದ-ರಿಂದ
ಬಳ-ಸು-ವುದು
ಬಳ್ಳಿಯ
ಬಸ್ತಿ
ಬಸ್ತಿ-ಯಂತಹ
ಬಹಳ
ಬಹಿಷ್ಕೃತ-ರಾದ-ವರು
ಬಾಗಿದ
ಬಾಗಿದಂತಿ-ರ-ಬೇಕು
ಬಾತು-ಕೋಳಿಯ
ಬಾಯಿಯ
ಬಾಯಿ-ಯನ್ನು
ಬಾಯಿ-ಯಲ್ಲಿ
ಬಾಯಿ-ಯೊಳಗೆ
ಬಾರ-ದಿದ್ದಾಗ
ಬಾರಿ
ಬಾಲಾ-ಲಿಂಗ-ನಾಸನ
ಬಾವು
ಬಾವು-ಟದ
ಬಿಗಿ-ಗೊಳಿ-ಸ-ಬೇಕು
ಬಿಗಿ-ಯಾಗಿ
ಬಿಗಿ-ಹಿಡಿದು
ಬಿಗಿ-ಹಿಡಿಯುವ
ಬಿಟ್ಟು
ಬಿಡ-ಬೇಕು
ಬಿಡು-ಗಡೆ
ಬಿಡು-ವಾಗ
ಬಿಡು-ವುದೇ
ಬಿದಿರಿನ
ಬಿದಿರಿ-ನಿಂದ
ಬಿದ್ದು
ಬಿಲ್ಲಿನ
ಬಿಲ್ಲಿ-ನಷ್ಟು
ಬಿಲ್ಲು-ಗಳಷ್ಟು
ಬಿಸಿ
ಬೀಜದ
ಬೀರ-ದಿದ್ದಾಗ
ಬೀರುವ
ಬೀಸುವ
ಬುಡಕ್ಕೆ
ಬುಡದ
ಬುದ್ಧಿಯ
ಬೂದಿಯ
ಬೂದಿ-ಯಿಂದ
ಬೂರುಗ
ಬೆಕ್ಕಿನ
ಬೆನ್ನ
ಬೆನ್ನನ್ನು
ಬೆನ್ನಿನ
ಬೆನ್ನು
ಬೆನ್ನೆಲು-ಬನ್ನು
ಬೆನ್ನೆಲುಬಿ-ನಿಂದ
ಬೆರಳು-ಗಳ
ಬೆರಳು-ಗಳನ್ನು
ಬೆರಳು-ಗಳಷ್ಟು
ಬೆರಳು-ಗಳಿಂದ
ಬೆರೆಸ-ಬೇಕು
ಬೆರೆಸಿ
ಬೆರೆ-ಸಿದ
ಬೆಳ-ಕಿಗೆ
ಬೆಳ-ಕಿನ
ಬೆಳ್ಳಿ
ಬೆವರಿ-ನೊಂದಿಗೆ
ಬೇಕಾ-ಗುತ್ತದೆ
ಬೇಕಾ-ಗುವ
ಬೇರೆ
ಬೇವಿನ
ಬೇಸಿ-ಗೆಯ
ಬೊಗಸೆ-ಗಳಷ್ಟು
ಬೋಧನೆ
ಬೋಧನೆಯ
ಬೋಧಿಸ-ಲಾ-ಗಿದೆ
ಬೋಧಿಸಿ-ದಂತೆಯೇ
ಬೋಧಿ-ಸಿದ್ದಾರೆ
ಬೋಳದ
ಬ್ರಹ್ಮಚರ್ಯ
ಬ್ರಹ್ಮಚರ್ಯವು
ಬ್ರಹ್ಮಚಾರಿ-ಗಳು
ಬ್ರಹ್ಮನ
ಬ್ರಹ್ಮ-ರಂಧ್ರ-ದ-ವರೆಗೆ
ಬ್ರಾಹ್ಮಣರು
ಬ್ರಾಹ್ಮಿ
ಭಂಗಿ
ಭಂಗಿ-ಯಲ್ಲಿ-ರ-ಬೇಕು
ಭಕ್ತಿ-ಯಿಂದ
ಭಗವಂತನ
ಭಗವತೇ
ಭಯ
ಭಯ-ವಿಲ್ಲದೇ
ಭವಾನಿ
ಭಸ್ತ್ರಾ
ಭಸ್ತ್ರಿಕಾ
ಭಾಗಕ್ಕೆ
ಭಾಗ-ಗಳನ್ನು
ಭಾಗದ
ಭಾಗ-ದಲ್ಲಿ
ಭಾಗ-ದಿಂದ
ಭಾಗ-ವನ್ನು
ಭಾಗವು
ಭಾರದ್ವಾಜಾ-ಸನ
ಭಾರ-ವನ್ನು
ಭಾರಾ-ಸನ
ಭೇದಿ-ಸಿದ
ಭೇದಿ-ಸುತ್ತದೆ
ಭೇದಿ-ಸುವ
ಭ್ರಾಮರಿ
ಭ್ರೂಮಧ್ಯದ
ಮಂಗನ
ಮಂಚದ
ಮಂಡೂಕ
ಮಂಡೂಕೋದರ
ಮಂತ್ರ
ಮಂತ್ರದ
ಮಂತ್ರ-ದೊಂದಿಗೆ
ಮಂತ್ರ-ವನ್ನು
ಮಂಥನ
ಮಂಥನ-ಪ್ರವೇಶ
ಮಕ್ಕಳು
ಮಗುಚಿ
ಮಗು-ವನ್ನು
ಮಟ್ಟಕ್ಕೆ
ಮಟ್ಟಿನ
ಮಡಚಿ
ಮಡಚಿದ
ಮಡಚುತ್ತಾ
ಮಡಚುವು-ದಕ್ಕೆ
ಮಣಿ-ಕಟ್ಟನ್ನು
ಮತ್ತು
ಮತ್ತೆ
ಮತ್ತೊಂದು
ಮತ್ಸ್ಯಾ-ಸನ
ಮತ್ಸ್ಯೇಂದ್ರ
ಮತ್ಸ್ಯೇಂದ್ರ-ನಾಥರು
ಮಥನ
ಮಧುರ
ಮಧ್ಯದ
ಮಧ್ಯ-ದಿಂದ
ಮಧ್ಯಮ
ಮಧ್ಯಾಹ್ನದ
ಮಧ್ಯೆ
ಮನಸಾ
ಮನಸ್ಸನ್ನು
ಮನಸ್ಸಿಗೆ
ಮನಸ್ಸಿ-ನಲ್ಲಿ
ಮನಸ್ಸಿ-ನಿಂದ
ಮನಸ್ಸು
ಮನುಷ್ಯನು
ಮಯೂರಾ-ಸನ
ಮಯೂರಾ-ಸನ-ದಲ್ಲಿ-ರು-ವಾಗ
ಮರ-ಗುಬ್ಬಚ್ಚಿಯ
ಮರಣ
ಮರದ
ಮರಳ-ಬೇಕು
ಮರು-ಕ್ಷಣವೇ
ಮರುಳು
ಮರ್ಮಾಸ್ಥಿ
ಮರ್ಮಾಸ್ಥಿಯ
ಮಲ
ಮಲಗ-ಬೇಕು
ಮಲಗಿ
ಮಲಗಿ-ಕೊಂಡು
ಮಲಗಿದ
ಮಲಗುತ್ತಾ
ಮಲ-ದ್ವಾರವು
ಮಲ-ವನ್ನು
ಮಲ-ವಿಸರ್ಜನೆ
ಮಲ-ವಿಸರ್ಜನೆಯ
ಮಲ್ಲಿಗೆ
ಮಲ್ಲಿಗೆಯ
ಮಸ್ತಕ-ದಲ್ಲಿ-ರುವ
ಮಹಾ-ಗಣ-ಪತಯೇ
ಮಹಾ-ಗಣ-ಪತಿ
ಮಹಾನ್
ಮಹಾ-ಬಂಧ
ಮಹಾ-ಬಂಧ-ದಲ್ಲಿ-ರುವಾ-ಗಲೇ
ಮಹಾ-ಬಂಧವು
ಮಹಾ-ಮುದ್ರೆ
ಮಹಾ-ಮುದ್ರೆಯು
ಮಹಾ-ರಾಷ್ಟ್ರ-ದಲ್ಲಿ
ಮಹಾ-ವೇಧ
ಮಹಾ-ವೇಧವು
ಮಹಿಳೆ-ಯರು
ಮಾಡ-ದಿದ್ದಾಗ
ಮಾಡ-ದಿರು-ವುದು
ಮಾಡ-ದಿರು-ವುದೇ
ಮಾಡದೆ
ಮಾಡ-ಬಹುದು
ಮಾಡ-ಬಾರದು
ಮಾಡ-ಬೇಕು
ಮಾಡ-ಲಾಗಿದೆ
ಮಾಡ-ಲಾಗು-ತ್ತದೆ
ಮಾಡಲು
ಮಾಡಲ್ಪಟ್ಟಿ-ರ-ಬೇಕು
ಮಾಡಿ
ಮಾಡಿ-ಕೊಳ್ಳ-ಬೇಕು
ಮಾಡಿದ
ಮಾಡಿ-ದರೆ
ಮಾಡಿ-ದಾಗ
ಮಾಡಿದ್ದು
ಮಾಡಿ-ಸು-ವುದು
ಮಾಡುತ್ತಾ
ಮಾಡುವ
ಮಾಡು-ವಂತೆ
ಮಾಡು-ವವರಿ-ಗಾಗಿ
ಮಾಡು-ವಾಗ
ಮಾಡು-ವುದ-ರಿಂದ
ಮಾಡು-ವುದಿಲ್ಲವೋ
ಮಾಡು-ವುದು
ಮಾಡು-ವುದೇ
ಮಾತನ್ನು
ಮಾತಿ-ನಿಂದ
ಮಾತು
ಮಾತೇ
ಮಾತ್ರ
ಮಾನಸಿಕ-ವಾಗಿ
ಮಾರ್ಗ-ದರ್ಶನ-ದಂತೆ
ಮಾರ್ಗ-ದರ್ಶನ-ದಲ್ಲಿ
ಮಾರ್ಗ-ದಿಂದ
ಮಾರ್ಗ-ವಾಗಿ
ಮಾರ್ಜಾರೋತ್ತಾನಾ-ಸನ
ಮಾರ್ಪಡು-ತ್ತದೆ
ಮಾಲಾ-ಸನ
ಮಿಲಿ
ಮಿಶ್ರಣ-ವನ್ನು
ಮೀನಿನ
ಮೀರಿದ
ಮುಂಗೈ-ಗಳನ್ನು
ಮುಂಗೈ-ಯಷ್ಟು
ಮುಂತಾದ
ಮುಂ-ದಕ್ಕೆ
ಮುಂದೆ
ಮುಂದೆಯೂ
ಮುಕ್ತ-ಗಳಿಸಿ
ಮುಕ್ತ-ನಾಗು-ತ್ತಾನೆ
ಮುಕ್ತ-ವಾಗಿ-ರು-ತ್ತವೆಯೋ
ಮುಕ್ತಾಯ
ಮುಕ್ತಾಯ-ವಾಯಿತು
ಮುಕ್ತಿ
ಮುಖ-ಗಳನ್ನು
ಮುಖ-ವನ್ನು
ಮುಚ್ಚಿ
ಮುಟ್ಟ-ಬೇಕು
ಮುಟ್ಟಿ-ಸ-ಬೇಕು
ಮುಟ್ಟು-ವಷ್ಟು
ಮುತ್ತಿನ
ಮುದ್ರೆ
ಮುದ್ರೆ-ಗಳನ್ನು
ಮುದ್ರೆ-ಗಳು
ಮುದ್ರೆಯ
ಮುದ್ರೆ-ಯನ್ನು
ಮುದ್ರೆ-ಯಿಂದ
ಮುದ್ರೆಯು
ಮುನಿಯು
ಮುಳ್ಳಿನ
ಮುಷ್ಟಿ
ಮುಷ್ಟಿ-ಗಳನ್ನು
ಮುಷ್ಟಿಯ
ಮುಷ್ಟಿ-ಯಿಂದ
ಮುಸಲಾ-ಸನ
ಮುಹೂರ್ತ-ದಲ್ಲಿ
ಮೂಗನ್ನು
ಮೂಗಿನ
ಮೂಗು
ಮೂತ್ರ
ಮೂತ್ರ-ಕೋಶದ
ಮೂತ್ರ-ಕೋಶದ-ವರೆಗೆ
ಮೂತ್ರ-ಕೋಶ-ವನ್ನು
ಮೂತ್ರ-ನಾಳಕ್ಕೆ
ಮೂತ್ರ-ನಾಳದ
ಮೂತ್ರ-ನಾಳ-ದಲ್ಲಿ
ಮೂತ್ರ-ನಾಳ-ದಲ್ಲಿ-ರುವ
ಮೂತ್ರ-ನಾಳ-ದೊಳಗೆ
ಮೂತ್ರ-ನಾಳವು
ಮೂತ್ರ-ವನ್ನು
ಮೂತ್ರ-ವಿಸರ್ಜನೆ
ಮೂತ್ರ-ವಿಸರ್ಜನೆ-ಯಾಗು-ತ್ತದೆ
ಮೂರು
ಮೂರ್ಚಾ
ಮೂರ್ಛಾ
ಮೂರ್ಛಾ-ವಸ್ಥೆಯು
ಮೂಲಕ
ಮೂಲಕವೇ
ಮೂಲದ
ಮೂಲ-ಬಂಧ
ಮೂಲಾ-ಧಾರ-ವನ್ನು
ಮೂಲೆ-ಗಳ-ವರೆಗೆ
ಮೂವತ್ತು
ಮೃತ್ಯು-ವನ್ನು
ಮೃದು
ಮೃದು-ವಾಗಿ
ಮೃದು-ವಾದ
ಮೆದು-ವಾದ
ಮೇಕೆಯ
ಮೇದಿನಿ
ಮೇಲಕ್ಕೆ
ಮೇಲಕ್ಕೆತ್ತಿ
ಮೇಲಿಂದ
ಮೇಲಿಟ್ಟ
ಮೇಲಿಟ್ಟು
ಮೇಲಿನ
ಮೇಲಿರಿ-ಸ-ಬೇಕು
ಮೇಲಿ-ರಿಸಿ
ಮೇಲಿ-ರುವ
ಮೇಲೂ
ಮೇಲೆ
ಮೇಲೆತ್ತ-ಬೇಕು
ಮೇಲೆತ್ತಿ
ಮೇಲೆತ್ತುತ್ತಾ
ಮೇಲೆ-ರುತ್ತದೆ
ಮೇಲೆ-ಳೆದು-ಕೊಳ್ಳ-ಬೇಕು
ಮೇಲೆ-ಳೆದು-ಕೊಳ್ಳುವ
ಮೇಲೆ-ಳೆದು-ಕೊಳ್ಳು-ವುದು
ಮೇಲೊತ್ತಿ
ಮೇಲ್ಭಾಗದ
ಮೇಲ್ಭಾಗ-ದಲ್ಲಿ
ಮೇಲ್ಭಾಗ-ದಲ್ಲಿ-ರುವ
ಮೇಲ್ಭಾಗ-ವನ್ನು
ಮೇಲ್ಮುಖ-ವಾಗಿ
ಮೇಲ್ಮುಖ-ವಾದ
ಮೇಷಾ-ಸನ
ಮೈಮೇಲೆ
ಮೊಗ್ಗಿ-ನಂತೆ
ಮೊಗ್ಗು-ಗ-ಳಂತೆ
ಮೊಣ-ಕಾಲನ್ನು
ಮೊಣ-ಕಾಲಿಗೆ
ಮೊಣ-ಕಾಲಿನ
ಮೊಣ-ಕಾಲು
ಮೊಣ-ಕಾಲು-ಗಳ
ಮೊಣ-ಕಾಲು-ಗಳನ್ನು
ಮೊಣ-ಕಾಲು-ಗಳಿಂದ
ಮೊಣ-ಕಾಲು-ಗಳಿಗೆ
ಮೊಣಕೈ
ಮೊಣ-ಕೈ-ಗಳನ್ನು
ಮೊದಲ
ಮೊದಲ-ನೆಯ-ದಾಗಿ
ಮೊದ-ಲಾದ
ಮೊದಲಿ-ಗಿಂತಲೂ
ಮೊದಲು
ಮೊಲದ
ಮೊಳ
ಮೊಳಕೆ-ಯಷ್ಟು
ಮೊಳೆಯ
ಮೊಸಳೆಯ
ಮೌನ
ಯಜ್ಞವೇ
ಯತ್ನಿಸ-ಬೇಕು
ಯಮ
ಯಮ-ಗಳನ್ನು
ಯಾರ
ಯಾರೂ
ಯಾವ
ಯಾವಾಗ
ಯಾವಾ-ಗಲೂ
ಯಾವು-ದಾ-ದರೂ
ಯಾವುದು
ಯಾವುದೇ
ಯುವತಿಯ
ಯೋಗ
ಯೋಗದ
ಯೋಗ-ದಲ್ಲಿ
ಯೋಗ-ಮುದ್ರಾಭ್ಯಾಸ-ಕ್ಕಾಗಿ
ಯೋಗ-ವನ್ನು
ಯೋಗಾ-ಭ್ಯಾಸಕ್ಕೆ
ಯೋಗಾ-ಭ್ಯಾಸದ
ಯೋಗಾ-ಭ್ಯಾಸ-ವನ್ನು
ಯೋಗಾ-ಭ್ಯಾಸವು
ಯೋಗಿ-ಗಳ
ಯೋಗಿ-ಗಳಿಗೆ
ಯೋಗಿಗೆ
ಯೋಗಿಯ
ಯೋಗಿ-ಯನ್ನು
ಯೋಗಿಯು
ಯೋಗೇಶ
ಯೋಚಿ-ಸ-ಬೇಕು
ಯೋಚಿ-ಸುತ್ತಾ
ಯೋನಿ-ಯಲ್ಲೇ
ಯೋನಿ-ಯೊಳ-ಗಿನ
ಯೋನಿ-ಯೊಳಗೆ
ಯೋನ್ಯಾ-ಸನ
ಯೌವನ-ಭರಿತರು
ರಂದ್ರ-ದಿಂದ
ರಂಧ್ರ-ಗಳ
ರಂಧ್ರ-ಗಳನ್ನು
ರಂಧ್ರ-ದೊಳಗೆ
ರಕ್ತ
ರಚಿಸ-ಲಾ-ಗಿದೆ
ರಚಿಸಿ-ದ್ದಾರೆ
ರಜಸ್ಸನ್ನು
ರತಿ
ರತ್ನ
ರತ್ನ-ವನ್ನು
ರಥಾ-ಸನ
ರಮಿಸ-ಬೇಕು
ರಸ
ರಾಜರು
ರೀತಿ
ರೀತಿಯ
ರೀತಿ-ಯಲ್ಲಿ
ರೀತಿ-ಯಾಗಿ
ರುಚಿಯ
ರುದ್ರನ
ರೂಪ
ರೂಪ-ವಾದ
ರೋಮ-ಗಳು
ರೋಮಾಂಚನ-ವಾಗುವ-ವರೆಗೆ
ಲಕ್ಷಣ
ಲಕ್ಷಣ-ಗಳು
ಲಕ್ಷಣ-ವಾ-ಗಿದೆ
ಲಜ್ಜೆ
ಲಭಿ-ಸುತ್ತದೆ
ಲಭಿ-ಸುವ
ಲವಂಗ
ಲವಣ
ಲವಣದ
ಲವಣ-ವನ್ನು
ಲವಾಲಾ
ಲಾಂಗಲಾ-ಸನ
ಲೀಟರ್
ಲುಂಠನಾ-ಸನ
ಲೇಪಿಸಿ-ರ-ಬೇಕು
ಲೈಂಗಿಕ
ಲೌಕಿಕ
ವಜ್ರ-ಕಂದರ
ವಜ್ರ-ದಿಂದ
ವಜ್ರಾ-ಸನ
ವಜ್ರಾ-ಸನ-ದಲ್ಲಿ
ವಜ್ರೊಲಿ
ವಜ್ರೋಲಿ
ವಜ್ರೋಲಿಗೆ
ವಜ್ರೋಲಿಯ
ವಜ್ರೋಲಿ-ಯನ್ನು
ವಜ್ರೋಲಿಯು
ವನ-ವಾಸಿ-ಗಳಿಗೆ
ವರಾಹಾ-ಸನ
ವರಾಹಾ-ಸನ-ವನ್ನು
ವರ್ಣನೆಯ
ವರ್ಣನೆ-ಯನ್ನು
ವರ್ಣನೆ-ಯಿಂದ
ವರ್ಣಿ-ಸುವ
ವಸ್ತು-ಗಳ
ವಸ್ತು-ಗಳನ್ನು
ವಸ್ತು-ಗಳಿಂದ
ವಾಂತಿ
ವಾಚಾ
ವಾತ್ಸಲ್ಯ-ವಿಲ್ಲದಿ-ರು-ವುದು
ವಾದ-ವಿವಾದ-ಗಳನ್ನು
ವಾನರಾ-ಸನ
ವಾನರಾ-ಸನ-ದಲ್ಲಿದ್ದು
ವಾಯು
ವಾಯು-ಗಳನ್ನು
ವಾಯು-ಗ್ರಹಣ
ವಾಯು-ವನ್ನು
ವಾಯು-ವಿನ
ವಾಯುವು
ವಾಸನೆ-ಯನ್ನು
ವಾಸು-ದೇವ
ವಾಸು-ದೇವನ
ವಾಸು-ದೇವಾಯ
ವಾಸ್ತವವೋ
ವಿಕಸಿತ-ಗೊಳ್ಳುತ್ತದೆ
ವಿಗ್ರಹ-ಗಳನ್ನು
ವಿಚಲನೆ-ಯಿಂದ
ವಿದ್ಯಾರ್ಥಿ-ಗಳು
ವಿಧದ
ವಿಧಾನ
ವಿಧಾನ-ಗಳ
ವಿಧಾನದ
ವಿಧಿಸ-ಲಾದ
ವಿಪರೀತ
ವಿಪರೀತ-ಕರಣೀ
ವಿಪರೀತ-ನೃತ್ಯಾ-ಸನ
ವಿಮಾನಾ-ಸನ
ವಿರತಾ-ಸನ
ವಿರೋಧಿ-ಗಳು
ವಿರೋಧಿಸು-ವವರ
ವಿಲೋಮ
ವಿವರ
ವಿವರಣೆ
ವಿವರ-ವಾಗಿ
ವಿವರಿಸ-ಲಾ-ಗಿದೆ
ವಿವಿಧ
ವಿಶಿಷ್ಟ-ವಾದ
ವಿಶೇಷ-ವಾಗಿ
ವಿಶ್ರಾಂತಿ
ವಿಶ್ರಾಂತಿ-ಗಾಗಿ
ವಿಷಯ
ವಿಷಯ-ಗಳು
ವಿಷಯ-ದಲ್ಲಿ
ವಿಷ್ಣು
ವಿಷ್ಣು-ವಿಗೆ
ವಿಷ್ಣು-ವಿನ
ವಿಸರ್ಜನೆ-ಗಾಗಿ
ವಿಸರ್ಜನೆ-ಯನ್ನು
ವಿಸರ್ಜಿ-ಸ-ಬೇಕು
ವಿಸ್ತರಿ-ಸುತ್ತದೆಯೋ
ವಿಹಿತ-ವಾ-ಗಿದೆ
ವೀಣೆ
ವೀರ್ಯ
ವೀರ್ಯಕ್ಕೆ
ವೀರ್ಯದ
ವೀರ್ಯ-ವನ್ನು
ವೀರ್ಯವು
ವೀರ್ಯ-ಸ್ಖಲನ
ವೃಂತಾ-ಸನ
ವೃಕಾ-ಸನ
ವೃತ್ತಾ-ಕಾರ-ದಲ್ಲಿ
ವೃದ್ಧಿ
ವೃದ್ಧಿ-ಯಾಗು-ತ್ತದೆ
ವೃಷಣ-ಗಳ
ವೃಷ-ಪಾದ-ಕ್ಷೇಪ-ದಿಂದ
ವೃಷ-ಪಾದ-ಕ್ಷೇಪಾ-ಸನ
ವೃಷ-ಪಾದ-ಕ್ಷೇಪಾ-ಸನ-ದಿಂದ
ವೃಷಭ
ವೇಗ-ವಾಗಿ
ವೇತ್ರಾ-ಸನ
ವೇತ್ರಾ-ಸನ-ದಲ್ಲಿ-ರುವಾ-ಗಲೇ
ವೇಳೆ
ವೈದ್ಯಕೀಯ
ವೈಮಾನಿಕ
ವೈರಾಗ್ಯ
ವೈಷ್ಣವ
ವ್ಯರ್ಥ
ವ್ಯರ್ಥ-ವಾ-ದಾಗ
ವ್ಯವಹಾರ-ಗಳಿಂದ
ವ್ಯಾಮೋಹ-ಕ್ಕೊಳಗಾದ-ವರು
ವ್ರತ-ಗಳಿಂದ
ಶಂಕ್ವಾ-ಸನ
ಶಕ್ತಿ
ಶಕ್ತಿ-ಚಾಲನ
ಶಕ್ತಿಯ
ಶಕ್ತಿ-ಯಿದ್ದರೂ
ಶಕ್ತಿಯೇ
ಶಬ್ದ
ಶಬ್ದ-ಗಳೊಂದಿಗೆ
ಶಬ್ದ-ದೊಂದಿಗೆ
ಶಬ್ದ-ವನ್ನು
ಶಮನ-ಗೊಳಿ-ಸಲು
ಶಮನ-ಗೊಳಿ-ಸು-ವುದು
ಶರೀರಕ್ಕೆ
ಶರೀರದ
ಶರೀರ-ದಲ್ಲಿ
ಶರೀರ-ವನ್ನು
ಶರೀರವು
ಶವ-ದಂತೆ
ಶವಾ-ಸನ
ಶಶಾ-ಸನ
ಶಸ್ತ್ರ
ಶಸ್ತ್ರ-ಗಳನ್ನು
ಶಸ್ತ್ರ-ಗಳಿಂದ
ಶಸ್ತ್ರ-ದಿಂದ
ಶಾಂತ-ವಾ-ದಾಗ
ಶಾಶ್ವತ
ಶಾಸ್ತ್ರ-ಗಳ
ಶಾಸ್ತ್ರ-ಗಳಲ್ಲಿ
ಶಾಸ್ತ್ರ-ಜ್ಞಾನ-ವಿ-ರುವ
ಶಾಸ್ತ್ರದ
ಶಿಕ್ಷಿ-ಸುವ
ಶಿಫಾರಸು
ಶಿರಸ್ಸಿನ
ಶಿವನ
ಶಿವನಿ
ಶಿವ-ನಿಗೆ
ಶಿವನಿ-ಯನ್ನು
ಶಿಶ್ನ
ಶಿಶ್ನದ
ಶಿಶ್ನ-ದೊಳಗೆ
ಶಿಶ್ನ-ದ್ವಾರದ
ಶಿಶ್ನ-ವನ್ನು
ಶಿಷ್ಯ-ರಿಗೆ
ಶಿಸ್ತಿನ
ಶೀತಲಿ
ಶೀರ್ಷಾ-ಸನದ
ಶುಂಠಿ
ಶುಕಾ-ಸನ
ಶುಕ್ತ್ಯಾ-ಸನ
ಶುಚಿ-ಗೊಳಿ-ಸಿದ
ಶುಚಿ-ಗೊಳಿ-ಸುವ
ಶುಚಿ-ಗೊಳಿ-ಸುವಿಕೆಯ
ಶುಚಿತ್ವ
ಶುಚಿ-ಯಾದ
ಶುದ್ಧ-ಗೊಂಡ-ವನಾಗು-ತ್ತಾನೆ
ಶುದ್ಧರು
ಶುದ್ಧಿ
ಶುದ್ಧಿ-ಕ್ರಿಯೆ-ಗಳನ್ನು
ಶುದ್ಧಿ-ಕ್ರಿಯೆ-ಗಳು
ಶುದ್ಧಿ-ಯಾಗುತ್ತದೆ
ಶುದ್ಧಿ-ಯಾಗುತ್ತವೆ
ಶುದ್ಧಿ-ಯಾದ
ಶುದ್ಧೀ-ಕರಣ
ಶುದ್ಧೀ-ಕರಿ-ಸ-ಬೇಕು
ಶುಭ-ದಾಯಕ-ವಾ-ಗಿದೆ
ಶುಭ-ದಾಯಕ-ವಾದ
ಶೂನ್ಯ
ಶೂಲಾ-ಸನ
ಶೋಧನಾ
ಶ್ಯೇನಾ-ಸನ
ಶ್ರದ್ಧೆ
ಶ್ರಮಿಸ-ಬೇಕು
ಶ್ರೀ
ಶ್ರೀಕೃಷ್ಣ
ಶ್ರೀಗಂಧ
ಶ್ರೀಗಂಧ-ವನ್ನು
ಶ್ರೀಗಣೇಶಾಯ
ಶ್ರೇಷ್ಠತೆ
ಶ್ಲಾಘಿಸು-ತ್ತಾರೆ
ಶ್ವೋತ್ತಾನಾ-ಸನ
ಶ್ವೋತ್ತಾನಾ-ಸನ-ದಲ್ಲಿ-ರುವಂತೆಯೇ
ಷಟ್ಕರ್ಮ
ಸಂಕುಚಿ-ಸ-ಬಲ್ಲನು
ಸಂಕುಚಿ-ಸ-ಬೇಕು
ಸಂಕುಚಿಸಿ
ಸಂಖ್ಯೆ-ಗಳನ್ನು
ಸಂಖ್ಯೆ-ಯನ್ನು
ಸಂಗ
ಸಂಗ್ರಹ-ಮಾಡ-ದಿರು-ವುದು
ಸಂಗ್ರಹ-ವಾಗುತ್ತದೆ
ಸಂಗ್ರಹ-ವಾದ
ಸಂಚಲನ-ವಾಗು-ವಂತೆ
ಸಂತೋಷ-ಪಡು-ವುದು
ಸಂದರ್ಭಕ್ಕೆ
ಸಂದರ್ಭ-ಗಳಲ್ಲೂ
ಸಂದಿ-ನಿಂದ
ಸಂಪತ್ತು
ಸಂಪರ್ಕ-ವನ್ನು
ಸಂಪೂರ್ಣ
ಸಂಪೂರ್ಣ-ವಾಗಿ
ಸಂಪೂರ್ಣ-ವಾಯಿತು
ಸಂಪ್ರದಾಯ
ಸಂಪ್ರದಾಯದ
ಸಂಬಂಧಿ-ಸಿದ್ದು
ಸಂಭೋಗ
ಸಂಭೋಗದ
ಸಂಭೋಗ-ದಲ್ಲಿ
ಸಂಭೋಗ-ವನ್ನು
ಸಂಭೋಗಿ-ಸ-ಬಹುದು
ಸಂಭೋಗಿಸಿ
ಸಂಭೋಗಿ-ಸುವ
ಸಂರಕ್ಷಿ-ಸಲು
ಸಂಶಯ-ವಿಲ್ಲ
ಸಂಸಾರ
ಸಂಸಾರದ
ಸಂಸ್ಕೃತ
ಸಕ್ಕರೆ
ಸಚ್ಚಿದಾನಂದ
ಸಡಿಲ-ಗೊಳಿ-ಸ-ಬೇಕು
ಸಡಿಲ-ಗೊಳಿ-ಸುತ್ತಾ
ಸಡಿಲ-ವಾಗು-ತ್ತದೆ
ಸಣ್ಣ
ಸಣ್ಣ-ದಾದ
ಸತ್ಯ
ಸತ್ಯ-ಪಾಲನೆ
ಸದ್ಗುರು-ವಿನ
ಸಮ-ತೋಲನ
ಸಮ-ತೋಲನ-ಗೊಳಿ-ಸ-ಬೇಕು
ಸಮ-ತೋಲನ-ದಲ್ಲಿ-ರ-ಬೇಕು
ಸಮಯದ
ಸಮಯ-ದಲ್ಲಿ
ಸಮಯ-ದಲ್ಲೂ
ಸಮಯ-ದ-ವರೆಗೆ
ಸಮರ್ಥ-ನಾಗು-ತ್ತಾನೆ
ಸಮಾಜ-ದಿಂದ
ಸಮಾನ-ವಾದ
ಸಮೀಪ-ವಿರ-ಬೇಕು
ಸಮೀಪ-ವಿರು-ತ್ತದೆ
ಸರಟಾ-ಸನ
ಸರದಿ-ಯಂತೆ
ಸರಿಯ-ಬೇಕು
ಸರಿ-ಯಾದ
ಸರ್ಪಾ-ಸನ
ಸರ್ವ-ವಾಯು
ಸಲಾಕೆ-ಗಳ
ಸಲಾಕೆ-ಗಳನ್ನು
ಸಲಾಕೆ-ಗಳಿಂದ
ಸಲಾಕೆಯ
ಸಲಾಕೆ-ಯನ್ನು
ಸಲು-ವಾಗಿ
ಸಹ
ಸಹಕರಿ-ಸು-ವುದು
ಸಹಕಾರಿ-ಯಾ-ಗಿದೆ
ಸಹಕಾರಿ-ಯಾದ
ಸಹಜ-ವಾಗಿ
ಸಹಜೋಲಿ
ಸಹನೀಯ-ವಾಗು-ತ್ತದೆ
ಸಹಾಯ
ಸಹಾಯ-ದಿಂದ
ಸಹಾಯ-ದಿಂ-ದಲೇ
ಸಹಾಯ-ವಿಲ್ಲದೆ
ಸಹಾಯ-ವಿಲ್ಲದೆಯೇ
ಸಹಿಸಿ-ಕೊಳ್ಳ-ಬೇಕು
ಸಹಿಸಿ-ಕೊಳ್ಳು-ವುದೇ
ಸಾಂದ್ರೀಕೃತ-ಗೊಂಡು
ಸಾಕಷ್ಟು
ಸಾಕ್ಷಾತ್ಕಾರದ
ಸಾಕ್ಷಾತ್ಕಾರ-ದಿಂದ
ಸಾಕ್ಷಾತ್ಕಾರ-ವನ್ನು
ಸಾಕ್ಷಾತ್ಕಾರ-ವಿಲ್ಲದೆ
ಸಾಧಕನು
ಸಾಧಕರು
ಸಾಧನ-ಗಳು
ಸಾಧನೆ
ಸಾಧನೆಯ
ಸಾಧನೆ-ಯಲ್ಲಿ
ಸಾಧ್ಯ-ವಾಗು-ತ್ತದೆ
ಸಾಧ್ಯ-ವಾಗು-ತ್ತದೆಯೋ
ಸಾಧ್ಯ-ವಾ-ದಾಗ-ಲೆಲ್ಲಾ
ಸಾಧ್ಯ-ವಿಲ್ಲ
ಸಾಮಗ್ರಿ-ಗಳ
ಸಾಮಥ್ರ್ಯ-ವನ್ನು
ಸಾಮರ್ಥ್ಯ
ಸಾಮರ್ಥ್ಯ-ವುಂಟಾಗುತ್ತದೆ
ಸಾರಿಸಿರ-ಬೇಕು
ಸಾಲ-ದಿದ್ದಾಗ
ಸಾವಿರ
ಸಾವಿರಾರು
ಸಾಸಿವೆ
ಸಿಕ್ಕಿದುದ-ರಲ್ಲೇ
ಸಿದ್ಧ
ಸಿದ್ಧ-ಪಡಿಸಿ-ಕೊಳ್ಳ-ಬೇಕು
ಸಿದ್ಧ-ಪುರುಷ-ನಾ-ದಾಗ
ಸಿದ್ಧ-ಯೋಗಿ-ಗಳ
ಸಿದ್ಧ-ವಾ-ದರೆ
ಸಿದ್ಧಿ
ಸಿದ್ಧಿಯ
ಸಿದ್ಧಿ-ಯಾಗುತ್ತದೆ
ಸಿದ್ಧಿ-ಯಾದ
ಸಿದ್ಧಿ-ಯಾ-ದಾಗ
ಸಿದ್ಧಿ-ಯಿಲ್ಲದೆ
ಸಿದ್ಧಿಯು
ಸಿದ್ಧಿ-ಯುಂಟಾಗುತ್ತದೆ
ಸಿದ್ಧಿಸಿ-ಕೊಳ್ಳಬಲ್ಲರು
ಸಿದ್ಧಿಸಿ-ಕೊಳ್ಳಲು
ಸಿದ್ಧಿಸಿ-ಕೊಳ್ಳುತ್ತಾರೆ
ಸಿದ್ಧಿಸಿ-ಕೊಳ್ಳುವ
ಸಿದ್ಧಿ-ಸಿದ
ಸಿದ್ಧಿಸಿ-ದಾಗ
ಸಿದ್ಧಿ-ಸಿದೆ
ಸಿದ್ಧಿ-ಸುತ್ತದೆ
ಸೀಟಿ
ಸೀತ್ಕಾರ
ಸೀಮನಿ
ಸುಂದರ
ಸುಂದರ-ವಾದ
ಸುಖಕ್ಕೆ
ಸುಖ-ಗಳನ್ನು
ಸುಖ-ವನ್ನು
ಸುಖಾ-ಸನ
ಸುಣ್ಣ-ದಿಂದ
ಸುತ್ತ-ಬೇಕು
ಸುತ್ತಿ
ಸುಶೀಲರು
ಸುಷುಮ್ನಾ
ಸೂಂ
ಸೂರ್ಯ
ಸೂರ್ಯ-ಭೇದನ
ಸೂರ್ಯ-ಭೇದನಂ
ಸೆಳೆ-ಯುವ
ಸೇರಿ-ದಂತೆ
ಸೇರಿಸ-ಬಾರದು
ಸೇರಿಸ-ಬೇಕು
ಸೇರಿಸಿ
ಸೇರಿ-ಸಿದ
ಸೇರಿ-ಸುತ್ತಿ-ರ-ಬೇಕು
ಸೇರಿ-ಸುವ
ಸೇರಿ-ಸು-ವುದು
ಸೇವನೆ
ಸೇವನೆ-ಯನ್ನು
ಸೇವಿ-ಸ-ಬೇಕು
ಸೇವಿ-ಸಿದ
ಸೈಂಧವ
ಸೊಂಟದ
ಸೋನಮಕ್ಕಿ
ಸೋಸಿ
ಸೋಸಿದ
ಸೌಂದರ್ಯ
ಸೌಂದರ್ಯದ
ಸೌತೆ-ಕಾಯಿ
ಸ್ಖಲನ
ಸ್ಖಲನಕ್ಕೆ
ಸ್ಖಲನ-ವಾಗುವಂತಹ
ಸ್ತಂಭನ
ಸ್ತನ-ಗಳನ್ನು
ಸ್ತನ-ಗಳು
ಸ್ತಬ್ಧ-ವಾಗು-ವುದೇ
ಸ್ತ್ರೀ
ಸ್ತ್ರೀಯ
ಸ್ತ್ರೀ-ಯನ್ನು
ಸ್ತ್ರೀಯರ
ಸ್ತ್ರೀಯರು
ಸ್ತ್ರೀಯ-ರೊಂದಿಗೆ
ಸ್ತ್ರೀಯು
ಸ್ತ್ರೀಯೊಂದಿಗೆ
ಸ್ಥಳ-ಗಳಲ್ಲೂ
ಸ್ಥಳ-ವಿರ-ಬೇಕು
ಸ್ಥಿತಿ-ಯಲ್ಲಿ
ಸ್ಥಿತಿ-ಯಲ್ಲಿದ್ದಾಗ
ಸ್ಥಿತಿ-ಯಲ್ಲೇ
ಸ್ಥಿರ-ಮನ-ಸ್ಥಿತಿ
ಸ್ಥಿರ-ವಾಗಿ
ಸ್ಥಿರ-ವಾಗಿ-ರಿಸಿ
ಸ್ಥಿರ-ವಾಗಿ-ರುವ
ಸ್ಥಿರ-ವಾದ
ಸ್ಥಿರ-ವಾ-ದಾಗ
ಸ್ನಾನ-ಕ್ಕಾಗಿ
ಸ್ನಾನಾದಿ-ಗಳಿಂದ
ಸ್ನಾಯು-ಗಳನ್ನು
ಸ್ನೇಹ-ಭಾವ-ವಿಲ್ಲ-ದಿ-ರುವುದೇ
ಸ್ನೇಹಿತನೂ
ಸ್ಪರ್ಶವೂ
ಸ್ಪರ್ಶಿ-ಸ-ಬೇಕು
ಸ್ಪರ್ಶಿ-ಸುವ
ಸ್ಪರ್ಶಿ-ಸು-ವುದು
ಸ್ಮರಣೆ-ಯನ್ನು
ಸ್ಮರಿಸ-ಬೇಕು
ಸ್ಮರಿಸುತ್ತಾ
ಸ್ವಂತ
ಸ್ವಚ್ಛ-ಗೊಳಿ-ಸ-ಬೇಕು
ಸ್ವಚ್ಛಗೊಳಿ-ಸುವ
ಸ್ವಧರ್ಮ-ವನ್ನು
ಸ್ವಭಾವ-ದ-ವರು
ಸ್ವರೂಪ-ವನ್ನು
ಸ್ವರ್ಗವು
ಸ್ವರ್ಗಾ-ರೋಹಣ
ಸ್ವರ್ಗಾ-ಸನ
ಸ್ವಲ್ಪ
ಸ್ವಲ್ಪ-ವಾಗಿ
ಸ್ವಲ್ಪ-ಸ್ವಲ್ಪ-ವಾಗಿ
ಸ್ವಸ್ತಿಕಾ-ಸನ
ಸ್ವಸ್ತಿಕಾ-ಸನ-ದಲ್ಲಿ
ಸ್ವೀಕರಿ-ಸ-ಬೇಕು
ಸ್ವೀಕರಿ-ಸುವು-ದನ್ನು
ಹಂತ
ಹಂತಕ್ಕೆ
ಹಂತದ-ವರೆಗೂ
ಹಂತದ-ವರೆಗೆ
ಹಂತ-ವನ್ನು
ಹಂತ-ವಾಗಿ
ಹಂದಿಯ
ಹಂಸಾ-ಸನ
ಹಗ್ಗ-ಗಳನ್ನು
ಹಗ್ಗದ
ಹಗ್ಗ-ವನ್ನು
ಹಠ-ಯೋಗ
ಹಠ-ಯೋಗದ
ಹಠ-ಯೋಗ-ದಂತಹ
ಹಠ-ಯೋಗಿಯು
ಹಣೆ-ಯನ್ನು
ಹತೋಟಿಗೆ
ಹತ್ತಿಯ
ಹತ್ತಿರ
ಹತ್ತಿರಕ್ಕೆ
ಹತ್ತಿರ-ದಲ್ಲಿ
ಹತ್ತಿರ-ವಿಟ್ಟು
ಹತ್ತು
ಹದ-ದಂತೆ
ಹದಿನಾರು
ಹದಿನೆಂಟು
ಹದಿನೈದು
ಹನಿಯನ್ನಷ್ಟೇ
ಹನಿ-ಯನ್ನು
ಹನ್ನೆರಡು
ಹರತಾಳ-ದಿಂದಲೂ
ಹರಳೆಣ್ಣೆ
ಹರಿಣಾ-ಸನ
ಹರಿತ-ಗೊಳಿಸಿ
ಹರಿತ-ವಾದ
ಹರಿತ-ಶರ
ಹರಿ-ಯುವ
ಹರಿವು
ಹಲ-ವಾರು
ಹಲ್ಲು-ಗಳನ್ನು
ಹಲ್ಲು-ಗಳಿಂದ
ಹಸಿ
ಹಸು-ವಿನ
ಹಸುವಿ-ನಿಂದ
ಹಾಕ-ಬೇಕು
ಹಾಕಿ-ಕೊಂಡು
ಹಾಕಿ-ಕೊಳ್ಳ-ಬೇಕು
ಹಾಕಿದ
ಹಾಗು
ಹಾಗೂ
ಹಾಗೆಯೇ
ಹಾದಿ-ಯಲ್ಲೇ
ಹಾದು-ಹೋಗು-ವಂತೆ
ಹಾನಿ
ಹಾರದ
ಹಾರಲು
ಹಾರಾ-ಟದ
ಹಾರಿಸ-ಬೇಕು
ಹಾರಿಸಿ
ಹಾಲನ್ನು
ಹಾಲಿ-ನಲ್ಲಿ
ಹಾಲು
ಹಾಲ್ಗರೆ-ಯು-ವುದು
ಹಾವಿನ
ಹಾಸಿಗೆಯ
ಹಾಸಿ-ರ-ಬೇಕು
ಹಿಂಜರಿಕೆ-ಯಿಲ್ಲದೆ
ಹಿಂಡುವ
ಹಿಂದಕ್ಕೆ
ಹಿಂದಿ-ನಿಂದ
ಹಿಂದೆ
ಹಿಂದೆಂದೂ
ಹಿಂಭಾಗ-ದಲ್ಲಿ
ಹಿಂಭಾಗ-ದಿಂದ
ಹಿಂಭಾಗ-ವನ್ನು
ಹಿಡಿದ
ಹಿಡಿದಿಟ್ಟು
ಹಿಡಿದಿಟ್ಟು-ಕೊಳ್ಳ-ಬೇಕು
ಹಿಡಿದಿಟ್ಟು-ಕೊಳ್ಳುವ
ಹಿಡಿದಿಡ-ಬೇಕು
ಹಿಡಿದು
ಹಿಡಿದು-ಕೊಂಡು
ಹಿಡಿದು-ಕೊಳ್ಳ-ಬೇಕು
ಹಿಡಿಯ-ಬೇಕು
ಹಿತಕರ-ವಾದ
ಹಿತ-ವನ್ನು
ಹಿಮ್ಮಡಿ
ಹಿಮ್ಮಡಿ-ಗಳನ್ನು
ಹಿಮ್ಮಡಿ-ಗಳಿಂದ
ಹಿಮ್ಮಡಿಯ
ಹಿಮ್ಮಡಿ-ಯನ್ನು
ಹಿಮ್ಮಡಿ-ಯಿಂದ
ಹಿಮ್ಮುಖ-ವಾಗಿ
ಹೀಗೆ
ಹೀರ-ಬೇಕು
ಹೀರಿ-ಕೊಳ್ಳುವ
ಹುಂ
ಹುಂಜದ
ಹುಣಸೆ
ಹುಬ್ಬು-ಗಳ
ಹುಲಿ
ಹುಲ್ಲಿನ
ಹುಲ್ಲು-ಗಳನ್ನು
ಹುಳು-ವಿನ
ಹೂವಿನ
ಹೃದಯದ
ಹೃದಯ-ದಲ್ಲಿ
ಹೃದಯ-ವನ್ನು
ಹೆಂಡತಿ
ಹೆಗಲ
ಹೆಗಲನ್ನು
ಹೆಗಲ-ವರೆಗೆ
ಹೆಚ್ಚಿನ
ಹೆಚ್ಚಿಸಿ-ಕೊಳ್ಳಲು
ಹೆಚ್ಚು
ಹೆಜ್ಜೆ-ಗಳ
ಹೆಬ್ಬೆರಳನ್ನು
ಹೆಬ್ಬೆರಳಿನ
ಹೆಬ್ಬೆರಳಿನಷ್ಟು
ಹೆಬ್ಬೆರಳಿ-ನಿಂದ
ಹೆಬ್ಬೆರಳು-ಗಳನ್ನು
ಹೆಸರಿಟ್ಟು
ಹೆಸರು-ಬೇಳೆ
ಹೇಗೆ
ಹೇಳಲಾ-ಗುತ್ತದೆ
ಹೇಳಿ-ಕೊಡ-ಲಾಗುತ್ತದೆ
ಹೇಳು-ವುದು
ಹೊಂದಿದ್ದಾನೆ
ಹೊಂದಿ-ರ-ಬೇಕು
ಹೊಂದಿ-ರುವ
ಹೊಂದುವ
ಹೊಕ್ಕಳಿನ
ಹೊಕ್ಕಳಿನ-ವರೆಗೆ
ಹೊಕ್ಕುಳ
ಹೊಕ್ಕುಳನ್ನು
ಹೊಕ್ಕುಳಿನ
ಹೊಟ್ಟೆ
ಹೊಟ್ಟೆಯ
ಹೊಟ್ಟೆ-ಯಂತೆ
ಹೊಟ್ಟೆ-ಯನ್ನು
ಹೊಟ್ಟೆ-ಯಲ್ಲಿ
ಹೊಟ್ಟೆ-ಯಲ್ಲಿನ
ಹೊಟ್ಟೆ-ಯಲ್ಲಿ-ರುವ
ಹೊಟ್ಟೆಯು
ಹೊಟ್ಟೆ-ಯೊಳಗೆ
ಹೊಡೆತ-ಗಳನ್ನು
ಹೊಡೆ-ದಂತಹ
ಹೊಡೆದು-ಕೊಳ್ಳ-ಬೇಕು
ಹೊಡೆ-ಯುವು-ದಕ್ಕೆ
ಹೊರಕ್ಕೆ
ಹೊರ-ಗಿನ
ಹೊರಗೆ
ಹೊರ-ತಾಗಿ
ಹೊರ-ತೆಗೆಯ-ಬೇಕು
ಹೊರ-ತೆಗೆಯುವ
ಹೊರ-ತೆಗೆಯು-ವುದು
ಹೊರ-ಬರುವ
ಹೊರ-ಬಿಡ-ಬೇಕು
ಹೊರ-ಬಿಡು-ವುದು
ಹೊರ-ಭಾಗ-ವನ್ನು
ಹೊರ-ಸೂಸು-ತ್ತದೆ
ಹೊರ-ಹಾಕ-ಬೇಕು
ಹೊರ-ಹಾಕ-ಲ್ಪಡುತ್ತವೆ
ಹೊರ-ಹಾಕುವ
ಹೊರ-ಹಾಕು-ವುದೇ
ಹೊರ-ಹೋಗು-ತ್ತದೆ
ಹೊಲಿಗೆ-ಗಳಂತಹ
ಹೊಳ್ಳೆ-ಗಳ
ಹೊಳ್ಳೆ-ಗಳು
ಹೊಳ್ಳೆಯ
ಹೊಸ್ತಿಲು
ಹೋಗ-ದಿದ್ದರೆ
ಹೋಗ-ಬೇಕು
ಹೋಗಿ
ಹೋದ
ಹೋಮ
}
